\begin{proof}[Proof of \cref{obj:thm:rectangular-upper}]
Fix the public parameters in the statement. Every constant constructed below depends only on these parameters.

\begin{enumerate}
\item Fix admissible \(n,m,h,J\) and \(P\in\mathcal M\). Define
\[
N=n+m,\qquad S=\alpha+\beta,\qquad
\delta=\frac hJ,\qquad
q=\binom{d+2}{2},\qquad k=qJ^d.
\]
In particular, \(N\ge n\ge2\) and \(\delta>0\).
By \cref{obj:lem:local-polynomial-projection-facts}, there is a constant
\(0<C_{\mathrm{loc}}<\infty\), uniform over these choices, and a coefficient vector
\(\vartheta\in\mathbb R^q\) such that
\[
p_\vartheta(x)=r(x)^\top\vartheta,\qquad
\sup_{x\in C_h}|\tau_P(x)-p_\vartheta(x)|
 \le C_{\mathrm{loc}}h^\gamma,\qquad
\|\vartheta\|_2\le C_{\mathrm{loc}},
\qquad
r_0^\top\vartheta=\tau_P(x_0).
\]
The same result supplies
\[
\|r_0\|_2\le C_{\mathrm{loc}},\qquad
\|e_Pr_u-\Pi_J(e_Pr_u)\|_{L_2(\nu_h)}
 \le C_{\mathrm{loc}}\delta^\alpha,\qquad
\|\mu_{0,P}-\Pi_J\mu_{0,P}\|_{L_2(\nu_h)}
 \le C_{\mathrm{loc}}\delta^\beta.
\]
Here \(x_0\in C_h\), so the bound on \(r_0\) follows from the stated uniform coarse-basis bound.

The finite Borel basis functions and the finite moment sums in
\cref{obj:def:moments} make both empirical moments finite Borel functions of the dataset.
The eigenvalue cutoff is Borel: its matrix superlevel set is the intersection of the closed unit-vector quadratic-form inequalities. Matrix inversion on the invertible branch and clipping are Borel. Thus the decision in
\cref{obj:def:estimator}, with value zero on the failed-cutoff branch, is finite and Borel on every dataset. Composing with the dataset coordinate gives a Borel decision on the original dataset and randomizer.%
% lean: CausalSmith.Stat.TwosamplePointcateAnnotationFrontier.tHat_measurable

\item Define the population moments and the localized Taylor remainder by
\[
\bar R=\mathbb E_{\mathsf E_{n,m}(P)}\widehat R_{h,J},
\qquad
Q_{P,h,J}=\mathbb E_{\mathsf E_{n,m}(P)}\widehat Q_{h,J},
\qquad
g_{\mathrm{rem}}:C_h\longrightarrow\mathbb R,\quad
g_{\mathrm{rem}}(x)=\tau_P(x)-p_\vartheta(x).
\]
All localized integrals below are over \(C_h\).

The observed-response identity and the exchangeability condition in
\cref{obj:def:primitive-class} give, almost surely,
\[
\mathbb E_P[AY\mid X]=e_P(X)\mu_{1,P}(X),
\qquad
\mathbb E_P[Y\mid X]
 =\mu_{0,P}(X)+e_P(X)\tau_P(X).
\]
Indeed, \(AY=AY(1)\), and conditional exchangeability factors its conditional mean. Applying the same argument to \((1-A)Y(0)\) and adding gives the second identity.

The role counts in \cref{obj:def:roles} satisfy
\[
\ell=\lfloor n/2\rfloor\ge1,\qquad
t=N-\ell\ge1.
\]
Their records are independent under the sampling law in
\cref{obj:ass:sample-law}. Consequently, taking expectations in
\cref{obj:def:moments}, using uniform design and factoring each rectangular product, yields
\[
\begin{aligned}
\bar R_u
&=\int r_u e_P\mu_{1,P}\,d\nu_h
 -\int e_Pr_u\,\Pi_J(\mu_{0,P}+e_P\tau_P)\,d\nu_h,\\
(Q_{P,h,J})_{uv}
&=\int r_u e_Pr_v\,d\nu_h
 -\int e_Pr_u\,\Pi_J(e_Pr_v)\,d\nu_h.
\end{aligned}
\]
The response in the rectangular correction is \(Y\), so its conditional mean is the second conditional mean displayed above.

For completeness, writing the orthonormal fine basis from
\cref{obj:def:projection} as
\[
\mathbf b_J=(b_{J,\iota})_{\iota=1}^k,
\]
its finite expansion gives, for arbitrary
\(f_1,f_2\in L_2(\nu_h)\),
\[
\int f_1\Pi_Jf_2\,d\nu_h
=
\sum_{\iota=1}^k
 \left(\int b_{J,\iota} f_1\,d\nu_h\right)
 \left(\int b_{J,\iota} f_2\,d\nu_h\right)
=
\int (\Pi_Jf_1)f_2\,d\nu_h.
\]
Orthonormality also gives
\[
\int f_1f_2\,d\nu_h-\int f_1\Pi_Jf_2\,d\nu_h
=
\int(f_1-\Pi_Jf_1)(f_2-\Pi_Jf_2)\,d\nu_h,
\]
and
\[
\|\Pi_Jf_1\|_{L_2(\nu_h)}^2
+\|f_1-\Pi_Jf_1\|_{L_2(\nu_h)}^2
=\|f_1\|_{L_2(\nu_h)}^2.
\]
These identities show that the population matrix displayed above is symmetric, so it agrees with the expectation of the symmetrized empirical matrix.

Multiplying its entries by \(\vartheta\), substituting
\(\mu_{1,P}=\mu_{0,P}+\tau_P\), and using the residual identity gives
\[
\begin{aligned}
\bar R_u-(Q_{P,h,J}\vartheta)_u
&=\int
 \bigl(e_Pr_u-\Pi_J(e_Pr_u)\bigr)
 \bigl(\mu_{0,P}-\Pi_J\mu_{0,P}\bigr)\,d\nu_h\\
&\quad+\int r_u e_Pg_{\mathrm{rem}}\,d\nu_h
 -\int \Pi_J(e_Pr_u)e_Pg_{\mathrm{rem}}\,d\nu_h.
\end{aligned}
\]
Cauchy--Schwarz and the two residual bounds from the preceding step imply
\[
\left|\int
 \bigl(e_Pr_u-\Pi_J(e_Pr_u)\bigr)
 \bigl(\mu_{0,P}-\Pi_J\mu_{0,P}\bigr)\,d\nu_h\right|
\le C_{\mathrm{loc}}^2\delta^S.
\]
The coarse coordinates are orthonormal, overlap bounds \(|e_P|\) by one, and projection contraction follows from the displayed Pythagorean identity. Hence
\[
\|r_u\|_{L_2(\nu_h)}=1,\qquad
\|e_Pr_u\|_{L_2(\nu_h)}\le1,\qquad
\|\Pi_J(e_Pr_u)\|_{L_2(\nu_h)}\le1,\qquad
\|e_Pg_{\mathrm{rem}}\|_{L_2(\nu_h)}
 \le C_{\mathrm{loc}}h^\gamma.
\]
Applying Cauchy--Schwarz to the remaining two integrals therefore gives
\[
|\bar R_u-(Q_{P,h,J}\vartheta)_u|
\le C_{\mathrm{loc}}^2\delta^S+2C_{\mathrm{loc}}h^\gamma
\le
(C_{\mathrm{loc}}^2+2C_{\mathrm{loc}})
(h^\gamma+\delta^S).
\]
Define
\[
C_{\mathrm{bias}}
=\sqrt q\,(C_{\mathrm{loc}}^2+2C_{\mathrm{loc}}).
\]
Squaring the coordinate bounds, summing over the \(q\) coordinates, and taking square roots proves
\[
\|\bar R-Q_{P,h,J}\vartheta\|_2
\le C_{\mathrm{bias}}(h^\gamma+\delta^S).
\]%
% lean: CausalSmith.Stat.TwosamplePointcateAnnotationFrontier.population_bias_coordinate_bound

\item Define the squared deviations, the bias scale, and the good event by
\[
\begin{aligned}
Z_R(\mathcal D_{n,m},U)
 &=\|\widehat R_{h,J}-\bar R\|_2^2,\\
Z_Q(\mathcal D_{n,m},U)
 &=\|\widehat Q_{h,J}-Q_{P,h,J}\|_F^2,\\
b_{\mathrm{bias}}&=h^\gamma+\delta^S,\\
g_\varepsilon&=\frac{\varepsilon(1-\varepsilon)}2,\\
G_{\mathrm{good}}
 &=\{(\mathcal D_{n,m},U):
       \lambda_{\min}(\widehat Q_{h,J})\ge g_\varepsilon\}.
\end{aligned}
\]
The first two functions ignore \(U\).

On \(G_{\mathrm{good}}\), the Rayleigh bound is
\[
\xi_{\mathrm{inv}}^\top\widehat Q_{h,J}\xi_{\mathrm{inv}}
 \ge g_\varepsilon\|\xi_{\mathrm{inv}}\|_2^2
\qquad(\xi_{\mathrm{inv}}\in\mathbb R^q).
\]
It excludes a nonzero kernel and thus makes the matrix invertible.
For arbitrary \(\xi_{\mathrm{rhs}}\in\mathbb R^q\), set
\[
\xi_{\mathrm{inv}}=\widehat Q_{h,J}^{-1}\xi_{\mathrm{rhs}}.
\]
Cauchy--Schwarz then gives
\[
g_\varepsilon\|\xi_{\mathrm{inv}}\|_2^2
\le \xi_{\mathrm{inv}}^\top\xi_{\mathrm{rhs}}
\le \|\xi_{\mathrm{inv}}\|_2\|\xi_{\mathrm{rhs}}\|_2,
\]
and consequently
\[
\|\widehat Q_{h,J}^{-1}\xi_{\mathrm{rhs}}\|_2
\le g_\varepsilon^{-1}\|\xi_{\mathrm{rhs}}\|_2.
\]
For \(\xi_{\mathrm{inv}}=0\) this is immediate; otherwise division by its positive norm proves it.

The inverse identity and the residual decomposition are
\[
\widehat Q_{h,J}^{-1}\widehat R_{h,J}-\vartheta
=
\widehat Q_{h,J}^{-1}
 (\widehat R_{h,J}-\widehat Q_{h,J}\vartheta),
\]
\[
\widehat R_{h,J}-\widehat Q_{h,J}\vartheta
=
(\widehat R_{h,J}-\bar R)
+(\bar R-Q_{P,h,J}\vartheta)
+(Q_{P,h,J}-\widehat Q_{h,J})\vartheta.
\]
The triangle inequality and rowwise Cauchy--Schwarz imply
\[
\|\widehat R_{h,J}-\widehat Q_{h,J}\vartheta\|_2
\le
\sqrt{Z_R}
+C_{\mathrm{bias}}b_{\mathrm{bias}}
+C_{\mathrm{loc}}\sqrt{Z_Q}.
\]
Because both arm means lie in \([0,1]\),
\[
\tau_P(x_0)\in[-1,1].
\]
Clipping toward this interval decreases absolute error: below its lower endpoint the clipped endpoint lies between the value and the target, above its upper endpoint the same holds, and inside the interval clipping preserves the value. Thus, on \(G_{\mathrm{good}}\),
\[
\begin{aligned}
|\widehat T_{h,J}-\tau_P(x_0)|
&\le g_\varepsilon^{-1}\|r_0\|_2
 \|\widehat R_{h,J}-\widehat Q_{h,J}\vartheta\|_2\\
&\le g_\varepsilon^{-1}C_{\mathrm{loc}}
 \left(\sqrt{Z_R}
 +C_{\mathrm{bias}}b_{\mathrm{bias}}
 +C_{\mathrm{loc}}\sqrt{Z_Q}\right).
\end{aligned}
\]
Define the positive finite constant
\[
C_{\mathrm{pt}}
=g_\varepsilon^{-1}C_{\mathrm{loc}}
 (1+C_{\mathrm{bias}}+C_{\mathrm{loc}}).
\]
All three coefficients in the preceding parentheses are at most
\(1+C_{\mathrm{bias}}+C_{\mathrm{loc}}\). On the complementary event the decision is zero and its error is at most \(1\). Combining these branches gives, on every sample,
\[
|\widehat T_{h,J}-\tau_P(x_0)|
\le
C_{\mathrm{pt}}
 (b_{\mathrm{bias}}+\sqrt{Z_R}+\sqrt{Z_Q})
+\mathbf1_{G_{\mathrm{good}}^c}.
\]%
% lean: CausalSmith.Stat.TwosamplePointcateAnnotationFrontier.rectangular_pointwise_error_bound

\item By \cref{obj:lem:population-positivity},
\[
v_{\mathrm{unit}}^\top Q_{P,h,J}v_{\mathrm{unit}}
\ge\varepsilon(1-\varepsilon)=2g_\varepsilon
\qquad
(v_{\mathrm{unit}}\in\mathbb R^q,\ \|v_{\mathrm{unit}}\|_2=1).
\]
For each such vector, finite Cauchy--Schwarz gives
\[
\left|
v_{\mathrm{unit}}^\top
(\widehat Q_{h,J}-Q_{P,h,J})v_{\mathrm{unit}}
\right|^2
\le Z_Q.
\]
If \(Z_Q<g_\varepsilon^2\), every unit-vector quadratic form of
\(\widehat Q_{h,J}\) is at least \(g_\varepsilon\), by subtracting this deviation from \(2g_\varepsilon\). Its infimum is therefore at least \(g_\varepsilon\). Consequently,
\[
G_{\mathrm{good}}^c
\ \Longrightarrow\
Z_Q\ge g_\varepsilon^2,
\qquad
\mathbf1_{G_{\mathrm{good}}^c}
\le g_\varepsilon^{-1}\sqrt{Z_Q}.
\]
Define
\[
\kappa_{\mathrm{mean}}=C_{\mathrm{pt}}+g_\varepsilon^{-1}.
\]
The preceding pointwise bound now becomes
\[
|\widehat T_{h,J}-\tau_P(x_0)|
\le
\kappa_{\mathrm{mean}}
 (b_{\mathrm{bias}}+\sqrt{Z_R}+\sqrt{Z_Q}).
\]

Define the variance scale by
\[
V_{\mathrm{var}}
=\frac1{nh^d}+\frac{qJ^d}{nNh^{2d}}.
\]
By \cref{obj:lem:rectangular-sampling-bound}, there is a uniform positive finite constant \(C_{\mathrm{samp}}\) such that
\[
\mathbb E_{\mathsf E_{n,m}(P)}Z_R
+\mathbb E_{\mathsf E_{n,m}(P)}Z_Q
\le C_{\mathrm{samp}}V_{\mathrm{var}}.
\]
Both deviations are nonnegative. Cauchy--Schwarz under the probability sampling law therefore gives
\[
\begin{aligned}
\mathbb E_{\mathsf E_{n,m}(P)}\sqrt{Z_R}
&\le
\sqrt{\mathbb E_{\mathsf E_{n,m}(P)}Z_R}
\le\sqrt{C_{\mathrm{samp}}}\sqrt{V_{\mathrm{var}}},\\
\mathbb E_{\mathsf E_{n,m}(P)}\sqrt{Z_Q}
&\le
\sqrt{\mathbb E_{\mathsf E_{n,m}(P)}Z_Q}
\le\sqrt{C_{\mathrm{samp}}}\sqrt{V_{\mathrm{var}}}.
\end{aligned}
\]
The decision and target both belong to \([-1,1]\), so their absolute difference is at most \(2\) and is integrable. Integrating the pointwise inequality and defining
\[
C_{\mathrm{mean}}
=\kappa_{\mathrm{mean}}(1+2\sqrt{C_{\mathrm{samp}}})
\]
proves
\[
\begin{aligned}
\mathcal R_{n,m}(\widehat T_{h,J},P)
&\le\kappa_{\mathrm{mean}}
 \left(b_{\mathrm{bias}}
       +2\sqrt{C_{\mathrm{samp}}}\sqrt{V_{\mathrm{var}}}\right)\\
&\le C_{\mathrm{mean}}
 (b_{\mathrm{bias}}+\sqrt{V_{\mathrm{var}}}).
\end{aligned}
\]%
% lean: CausalSmith.Stat.TwosamplePointcateAnnotationFrontier.rectangular_risk_sqrt_variance

\item Define
\[
\rho_{\mathrm{lab}}=(nh^d)^{-1/2},\qquad
\rho_{\mathrm{rect}}=(nNh^d\delta^d)^{-1/2},\qquad
a_{\mathrm{noise}}=\rho_{\mathrm{lab}}+\rho_{\mathrm{rect}},
\qquad
\kappa_{\mathrm{rank}}=1+\sqrt q.
\]
All these quantities are nonnegative, and
\(\kappa_{\mathrm{rank}}\ge1\). Since \(\delta=h/J\),
\[
\rho_{\mathrm{lab}}^2=\frac1{nh^d},\qquad
\rho_{\mathrm{rect}}^2
=\frac{J^d}{nNh^{2d}},\qquad
V_{\mathrm{var}}
=\rho_{\mathrm{lab}}^2+q\rho_{\mathrm{rect}}^2.
\]
It follows that
\[
V_{\mathrm{var}}
\le
(\rho_{\mathrm{lab}}+\sqrt q\,\rho_{\mathrm{rect}})^2,
\qquad
\rho_{\mathrm{lab}}+\sqrt q\,\rho_{\mathrm{rect}}
\le\kappa_{\mathrm{rank}}a_{\mathrm{noise}},
\]
and hence
\[
\sqrt{V_{\mathrm{var}}}
\le\kappa_{\mathrm{rank}}a_{\mathrm{noise}}.
\]
Define
\[
C_{\mathrm{fixed}}
=C_{\mathrm{mean}}\kappa_{\mathrm{rank}}+2.
\]
Using also
\(b_{\mathrm{bias}}\le\kappa_{\mathrm{rank}}b_{\mathrm{bias}}\), the preceding risk bound gives
\[
\mathcal R_{n,m}(\widehat T_{h,J},P)
\le C_{\mathrm{fixed}}
 (b_{\mathrm{bias}}+a_{\mathrm{noise}}).
\]
If \(b_{\mathrm{bias}}+a_{\mathrm{noise}}\le1\), this is the desired bound with the minimum. If
\(b_{\mathrm{bias}}+a_{\mathrm{noise}}>1\), the pointwise diameter bound gives
\[
\mathcal R_{n,m}(\widehat T_{h,J},P)
\le2\le C_{\mathrm{fixed}}.
\]
Thus, uniformly over \(P\in\mathcal M\),
\[
\mathcal R_{n,m}(\widehat T_{h,J},P)
\le C_{\mathrm{fixed}}
\min\left\{1,\,
h^\gamma+\delta^S+(nh^d)^{-1/2}
 +(nNh^d\delta^d)^{-1/2}\right\}.
\]%
% lean: CausalSmith.Stat.TwosamplePointcateAnnotationFrontier.rectangular_fixed_tuning_risk

\item We next establish the public tuning bound. Define
\[
\Delta_{\mathrm{o}}=2\gamma+d,\qquad
\Delta=2\gamma+d+\frac{\gamma d}{S},
\qquad
c_{\mathrm{lab}}=2^{\Delta_{\mathrm{o}}/2},
\qquad
c_{\mathrm{rect}}
=2^{d/2}2^{\Delta/2}+2^{\Delta_{\mathrm{o}}},
\qquad
C_{\mathrm{bal}}=2+c_{\mathrm{lab}}+c_{\mathrm{rect}}.
\]
The parameter domain gives \(S>0\), \(\gamma>0\),
\(\Delta_{\mathrm{o}}>0\), and \(\Delta>0\), so these constants are positive and finite.

For fixed \(n\ge2\) and \(m\ge0\), define the two bandwidth orders
\[
u_{\mathrm{o}}=n^{-1/\Delta_{\mathrm{o}}},
\qquad
u_{\mathrm{rect}}=(nN)^{-1/\Delta}.
\]
Because \(n\ge1\) and \(nN\ge1\),
\[
0<u_{\mathrm{o}}\le1,\qquad
0<u_{\mathrm{rect}}\le1.
\]
The definition in \cref{obj:def:tuning} consequently gives
\[
h_*=\frac12\max\{u_{\mathrm{o}},u_{\mathrm{rect}}\},
\qquad
0<h_*\le\frac12,\qquad
u_{\mathrm{o}}\le2h_*,
\qquad
u_{\mathrm{rect}}\le2h_*.
\]
Moreover,
\[
\begin{aligned}
h_*^\gamma
&\le \max\{u_{\mathrm{o}}^\gamma,u_{\mathrm{rect}}^\gamma\}\\
&=\max\left\{
 n^{-\gamma/\Delta_{\mathrm{o}}},
 (nN)^{-\gamma/\Delta}\right\}
=r_{\mathrm{up}}(n,m),
\end{aligned}
\]
using \cref{obj:def:sharp-rate}.

Taking logarithms of \(u_{\mathrm{o}}\le2h_*\) gives
\[
-\log n
\le\Delta_{\mathrm{o}}(\log2+\log h_*).
\]
Therefore
\[
\begin{aligned}
\log\bigl((nh_*^d)^{-1/2}\bigr)
&=-\frac12(\log n+d\log h_*)\\
&\le\frac{\Delta_{\mathrm{o}}}{2}\log2
   +\gamma\log h_*,
\end{aligned}
\]
where \(\Delta_{\mathrm{o}}-d=2\gamma\). Exponentiating yields
\[
(nh_*^d)^{-1/2}\le c_{\mathrm{lab}}h_*^\gamma.
\]%
% lean: CausalSmith.Stat.TwosamplePointcateAnnotationFrontier.bandwidth_monomial_balance

\item Suppose first that \(S<\gamma\). Define
\[
z_{\mathrm{grid}}=h_*^{-(\gamma/S-1)}.
\]
Then \(z_{\mathrm{grid}}\ge1\), and \cref{obj:def:tuning} gives
\[
J_*=\lceil z_{\mathrm{grid}}\rceil,\qquad
1\le J_*,\qquad
z_{\mathrm{grid}}\le J_*<z_{\mathrm{grid}}+1
\le2z_{\mathrm{grid}}.
\]
Since \(h_*/z_{\mathrm{grid}}=h_*^{\gamma/S}\), these inequalities imply
\[
0<\frac{h_*}{J_*},\qquad
\frac12h_*^{\gamma/S}
\le\frac{h_*}{J_*}\le h_*^{\gamma/S},
\qquad
\left(\frac{h_*}{J_*}\right)^S\le h_*^\gamma.
\]
The lower cell-side bound gives
\[
\log(h_*/J_*)\ge\frac{\gamma}{S}\log h_*-\log2.
\]
Substituting this inequality into the logarithm of the rectangular sampling term and exponentiating proves
\[
\left(nNh_*^d(h_*/J_*)^d\right)^{-1/2}
\le
2^{d/2}
\left(nNh_*^{d+\gamma d/S}\right)^{-1/2}.
\]
Similarly, \(u_{\mathrm{rect}}\le2h_*\) gives
\[
-\log(nN)\le\Delta(\log2+\log h_*).
\]
As \(\Delta-(d+\gamma d/S)=2\gamma\),
\[
\begin{aligned}
\log\left(\left(nNh_*^{d+\gamma d/S}\right)^{-1/2}\right)
&=-\frac12\left(\log(nN)
 +(d+\gamma d/S)\log h_*\right)\\
&\le\frac{\Delta}{2}\log2+\gamma\log h_*.
\end{aligned}
\]
Consequently,
\[
\left(nNh_*^d(h_*/J_*)^d\right)^{-1/2}
\le2^{d/2}2^{\Delta/2}h_*^\gamma
\le c_{\mathrm{rect}}h_*^\gamma.
\]%
% lean: CausalSmith.Stat.TwosamplePointcateAnnotationFrontier.ceil_grid_side_bounds

\item Suppose instead that \(S\ge\gamma\), including equality. The exponent in
\cref{obj:def:tuning} is then zero, so
\[
J_*=1,\qquad
0<\frac{h_*}{J_*}=h_*,
\qquad
\left(\frac{h_*}{J_*}\right)^S
=h_*^S\le h_*^\gamma,
\]
because \(0<h_*\le1\).

The oracle bandwidth inequality also reads
\[
(n^2)^{-1/(2\Delta_{\mathrm{o}})}
=u_{\mathrm{o}}\le2h_*.
\]
Taking logarithms gives
\[
-\log(n^2)
\le2\Delta_{\mathrm{o}}(\log2+\log h_*).
\]
Using \(2\Delta_{\mathrm{o}}-2d=4\gamma\), we obtain
\[
\begin{aligned}
\log\left((n^2h_*^{2d})^{-1/2}\right)
&=-\frac12\left(\log(n^2)+2d\log h_*\right)\\
&\le\Delta_{\mathrm{o}}\log2+2\gamma\log h_*.
\end{aligned}
\]
Since \(N\ge n\) and \(h_*^{2\gamma}\le h_*^\gamma\),
\[
\begin{aligned}
\left(nNh_*^d(h_*/J_*)^d\right)^{-1/2}
&=(nNh_*^{2d})^{-1/2}\\
&\le(n^2h_*^{2d})^{-1/2}\\
&\le2^{\Delta_{\mathrm{o}}}h_*^{2\gamma}\\
&\le c_{\mathrm{rect}}h_*^\gamma.
\end{aligned}
\]
Thus both cases give admissible tuning and the four-term balance
\[
\begin{gathered}
0<h_*\le\frac12,\qquad 1\le J_*,\\
h_*^\gamma+(h_*/J_*)^S+(nh_*^d)^{-1/2}
+\left(nNh_*^d(h_*/J_*)^d\right)^{-1/2}\\
\le(2+c_{\mathrm{lab}}+c_{\mathrm{rect}})h_*^\gamma
\le C_{\mathrm{bal}}r_{\mathrm{up}}(n,m).
\end{gathered}
\]%
% lean: CausalSmith.Stat.TwosamplePointcateAnnotationFrontier.tuning_balance

\item Finally, define one constant for both assertions:
\[
C=C_{\mathrm{fixed}}(C_{\mathrm{bal}}+1).
\]
It is positive and finite, and
\[
C_{\mathrm{fixed}}\le C,\qquad
C_{\mathrm{fixed}}C_{\mathrm{bal}}\le C.
\]
The fixed-tuning envelope has a nonnegative minimum, so enlarging its constant from \(C_{\mathrm{fixed}}\) to \(C\) and taking the supremum over \(P\in\mathcal M\) proves the first risk assertion.

For the second assertion, the tuning bounds make
\((h_*,J_*)\) admissible for the fixed-tuning result. The identity in
\cref{obj:def:tuning} supplies
\[
\widehat T_{\mathrm{up}}=\widehat T_{h_*,J_*},
\]
and hence also its finiteness and Borel measurability. Using the minimum's upper bound by its four-term argument gives, for every \(P\in\mathcal M\),
\[
\begin{aligned}
\mathcal R_{n,m}(\widehat T_{\mathrm{up}},P)
&\le C_{\mathrm{fixed}}
 \left\{h_*^\gamma+(h_*/J_*)^S+(nh_*^d)^{-1/2}
 +\left(nNh_*^d(h_*/J_*)^d\right)^{-1/2}\right\}\\
&\le C_{\mathrm{fixed}}C_{\mathrm{bal}}r_{\mathrm{up}}(n,m)\\
&\le Cr_{\mathrm{up}}(n,m).
\end{aligned}
\]
The rate is nonnegative as the maximum of two positive powers of positive sample counts. All constants are uniform over the primitive class, so taking the supremum proves the tuned assertion.%
% lean: CausalSmith.Stat.TwosamplePointcateAnnotationFrontier.rectangular_upper
\end{enumerate}
\end{proof}

\begin{proof}[Proof of \cref{obj:thm:sharp-annotation-frontier}]
Fix the public parameters as in the statement. Throughout the proof, use
\[
N=n+m,\qquad S=\alpha+\beta,\qquad
\Delta=2\gamma+d+\frac{\gamma d}{S},\qquad
S_{\mathrm{crit}}=\frac{\gamma d}{2\gamma+d},\qquad
q_*=\frac{\gamma d}{S(2\gamma+d)},
\]
and
\[
r_o(n)=n^{-\gamma/(2\gamma+d)},\qquad
r_*=\max\{r_o(n),(nN)^{-\gamma/\Delta}\},
\]
where the last equality is \cref{obj:def:sharp-rate}. All sample sizes below satisfy \(n,m\in\mathbb N\) and \(n\ge2\).

\begin{enumerate}
\item \emph{Supplied-propensity lower bound.}
By \cref{obj:thm:supplied-propensity}, there is a constant \(c_O>0\), depending on the fixed public parameters, such that
\[
c_O r_o(n)\le R^e(n,m)
\qquad\text{for every }n,m.
\]
Every admissible decision \(T\) induces a supplied-propensity decision by
\[
T^e(\mathcal D_{n,m},U,e)=T(\mathcal D_{n,m},U)
\]
for every function \(e:[0,1]^d\to\mathbb R\). For each supplied function this decision is Borel, and its loss under each \(P\in\mathcal M\) is
\[
\int
\left|T^e(\mathcal D_{n,m},U,e_P)-\tau_P(x_0)\right|
\,d\mathsf E_{n,m}(P)
=
\mathcal R_{n,m}(T,P).
\]
Taking the supremum over the same primitive class and then the infimum over decisions, using \cref{obj:def:oracle-risk,obj:def:minimax}, gives
\[
R^e(n,m)\le R(n,m).
\]
Consequently,
\[
c_O r_o(n)\le R(n,m)
\qquad\text{for every }n,m.
\]
% lean: CausalSmith.Stat.TwosamplePointcateAnnotationFrontier.oracleRisk_le_minimaxRisk

\item \emph{A uniform lower constant.}
First suppose \(S<S_{\mathrm{crit}}\). By
\cref{obj:lem:nuisance-frontier-converse}, there is a constant \(c_B>0\), depending on the fixed public parameters, such that
\[
c_B(nN)^{-\gamma/\Delta}\le R(n,m)
\qquad\text{whenever }n\le N\le n^{q_*}.
\]
Choose
\[
c=\min\{c_O,c_B\}>0.
\]

We establish the rate comparisons used in this case and in the remaining smoothness case. The parameter conditions imply
\[
S>0,\qquad \gamma>0,\qquad 2\gamma+d>0,\qquad
\Delta>0,\qquad q_*>0,\qquad N\ge n\ge2.
\]
Direct substitution gives
\[
q_*=
\frac{\gamma d}{S(2\gamma+d)}
=
\frac{S_{\mathrm{crit}}}{S},
\]
and
\[
\Delta
=
(2\gamma+d)+\frac{\gamma d}{S}
=
(2\gamma+d)
\left(1+\frac{\gamma d}{S(2\gamma+d)}\right)
=
(2\gamma+d)(1+q_*).
\]
% lean: CausalSmith.Stat.TwosamplePointcateAnnotationFrontier.rate_exponent_identities

Since \(n>0\), the laws of real powers therefore give the boundary identity
\[
\begin{aligned}
(n\,n^{q_*})^{-\gamma/\Delta}
&=(n^{1+q_*})^{-\gamma/\Delta}\\
&=n^{-(1+q_*)\gamma/\Delta}\\
&=n^{-\gamma/(2\gamma+d)}
=r_o(n).
\end{aligned}
\]
The exponent \(-\gamma/\Delta\) is negative. Thus, when \(N\ge n^{q_*}\),
\[
0<n\,n^{q_*}\le nN
\quad\Longrightarrow\quad
(nN)^{-\gamma/\Delta}
\le(n\,n^{q_*})^{-\gamma/\Delta}
=r_o(n).
\]
When \(N\le n^{q_*}\), instead,
\[
0<nN\le n\,n^{q_*}
\quad\Longrightarrow\quad
r_o(n)
=(n\,n^{q_*})^{-\gamma/\Delta}
\le(nN)^{-\gamma/\Delta}.
\]
For completeness, if \(S\ge S_{\mathrm{crit}}\), then
\[
q_*=\frac{S_{\mathrm{crit}}}{S}\le1,
\qquad
n^{q_*}\le n^1=n\le N,
\]
where the power comparison uses \(n\ge1\). Taking the maximum in the definition of \(r_*\) consequently gives
\[
\begin{aligned}
S\ge S_{\mathrm{crit}}
&\quad\Longrightarrow\quad r_*=r_o(n),\\
S<S_{\mathrm{crit}},\ N\le n^{q_*}
&\quad\Longrightarrow\quad r_*=(nN)^{-\gamma/\Delta},\\
S<S_{\mathrm{crit}},\ N\ge n^{q_*}
&\quad\Longrightarrow\quad r_*=r_o(n).
\end{aligned}
\]
The boundary identity also gives
\[
N=n^{q_*}
\quad\Longrightarrow\quad
(nN)^{-\gamma/\Delta}=r_o(n).
\]
% lean: CausalSmith.Stat.TwosamplePointcateAnnotationFrontier.rate_branches

Return to the case \(S<S_{\mathrm{crit}}\). If \(N\le n^{q_*}\), then \(N\ge n\), so the nuisance lower bound applies and yields
\[
c r_*
=c(nN)^{-\gamma/\Delta}
\le c_B(nN)^{-\gamma/\Delta}
\le R(n,m).
\]
If \(N>n^{q_*}\), the supplied-propensity lower bound yields
\[
c r_*=c r_o(n)\le c_O r_o(n)\le R(n,m).
\]
Here the comparisons of constants preserve the inequalities because both rate terms are nonnegative.

In the remaining case \(S\ge S_{\mathrm{crit}}\), choose
\[
c=c_O>0.
\]
Then \(r_*=r_o(n)\), and hence
\[
c r_*=c_O r_o(n)\le R(n,m).
\]
This constructs a positive constant \(c\), depending on the fixed public parameters and valid simultaneously for every admissible sample-size pair, such that
\[
c r_*\le R(n,m).
\]
% lean: CausalSmith.Stat.TwosamplePointcateAnnotationFrontier.sharpRate_minimax_lower

\item \emph{Attainment and the upper constant.}
By \cref{obj:thm:rectangular-upper}, there is a real constant \(C_U>0\), depending on the fixed public parameters, such that the publicly tuned decision is Borel measurable and its risk under every \(P\in\mathcal M\) satisfies
\[
\mathcal R_{n,m}(\widehat T_{\mathrm{up}},P)
\le C_U r_{\mathrm{up}}(n,m).
\]
By \cref{obj:def:sharp-decision,obj:def:sharp-rate},
\[
T_*=\widehat T_{\mathrm{up}},
\qquad
r_{\mathrm{up}}(n,m)=r_*.
\]
Set
\[
C=\max\{c,C_U\}.
\]
Then \(0<c\le C<\infty\). For every admissible sample-size pair, \(T_*\) is Borel measurable, and taking the supremum of the pointwise risk bound gives
\[
\sup_{P\in\mathcal M}\mathcal R_{n,m}(T_*,P)
\le C_U r_*
\le C r_*.
\]
The last inequality follows from \(C_U\le C\) and \(r_*\ge0\).
% lean: CausalSmith.Stat.TwosamplePointcateAnnotationFrontier.sharp_annotation_frontier

\item \emph{Minimax comparison and arbitrary decisions.}
Since \(T_*\) belongs to the decision class in \cref{obj:def:minimax},
\[
R(n,m)\le
\sup_{P\in\mathcal M}\mathcal R_{n,m}(T_*,P).
\]
Together with the bounds already established, this proves
\[
c r_*
\le R(n,m)
\le\sup_{P\in\mathcal M}\mathcal R_{n,m}(T_*,P)
\le C r_*.
\]
For any Borel real-valued decision \(T\) using the original dataset and the independent uniform randomizer, the same infimum comparison gives
\[
c r_*\le R(n,m)
\le\sup_{P\in\mathcal M}\mathcal R_{n,m}(T,P).
\]
The rate comparisons and boundary identity established above supply all three stated branches and their equality at \(N=n^{q_*}\).
% lean: CausalSmith.Stat.TwosamplePointcateAnnotationFrontier.sharp_annotation_frontier
\end{enumerate}
\end{proof}

\begin{proof}[Proof of \cref{obj:thm:supplied-propensity}]
Fix the public parameters in the stated domain. Throughout, inverse-propensity quotients are assigned value zero when their denominator is zero:
\[
\frac{t_{\mathrm{num}}}{0}:=0
\qquad(t_{\mathrm{num}}\in\mathbb R).
\]
The designated propensity is extended by zero outside \([0,1]^d\).

\begin{enumerate}
\item Define the constant-nuisance subfamily and its supplied-propensity risk by
\[
\mathcal M_{1/2}
=
\left\{
P\in\mathcal M:
e_P(x)=\mu_{0,P}(x)=\frac12
\text{ for every }x\in[0,1]^d
\right\},
\]
\[
\mathcal R^e_{n,m}(T^e,P)
=
\int
\left|T^e(\mathcal D_{n,m},U,e_P)-\tau_P(x_0)\right|
\,d\mathsf E_{n,m}(P),
\qquad
R^e_{1/2}(n,m)
=
\inf_{T^e}\sup_{P\in\mathcal M_{1/2}}
\mathcal R^e_{n,m}(T^e,P).
\]
These loss integrals may take the value \(+\infty\). Since the subfamily uses the same decision class as \cref{obj:def:oracle-risk}, restriction of the supremum gives
\[
R^e_{1/2}(n,m)\le R^e(n,m).
\]
It therefore suffices to establish a lower bound for \(R^e_{1/2}(n,m)\), a uniform upper bound for the displayed smoother, and its admissibility. We first construct the lower-bound family.
% lean: CausalSmith.Stat.TwosamplePointcateAnnotationFrontier.oracleSubfamilyRisk_le_oracleRisk

\item Use the smooth profile from \cref{obj:def:marked-handle}, with the explicit formulas
\[
\chi(\upsilon)=
\begin{cases}
\exp(-1/\upsilon),&\upsilon>0,\\
0,&\upsilon\le0,
\end{cases}
\qquad \upsilon\in\mathbb R,
\qquad
B(w)=\prod_{\iota_{\mathrm{cov}}=1}^d\frac{\chi(1-4w_{\iota_{\mathrm{cov}}}^2)}{\chi(1)},
\qquad w\in\mathbb R^d.
\]
For this argument, extend the scaled-bump formula to every positive localization width and every covariate vector:
\[
B_h(x)=B\bigl((x-x_0)/h\bigr),
\qquad
C_h=\prod_{j=1}^d
\left[\frac12-\frac h2,\frac12+\frac h2\right],
\qquad h>0,\quad x\in\mathbb R^d.
\]
Smoothness of the exponential glue function and finite products give
\[
B\in C^\infty(\mathbb R^d),
\qquad
\operatorname{supp}(B)\subseteq[-1/2,1/2]^d,
\qquad
0\le B\le1,
\qquad
B(0)=1.
\]
In particular,
\[
0\le B_h\le1,\qquad
B_h(x)=0\quad(x\notin C_h),\qquad
B_h(x_0)=1.
\]

We need a uniform Hölder bound after scaling. Define
\[
p_*=\lceil\gamma\rceil-1,
\qquad
s_o=\gamma-p_*,
\qquad
0\le p_*\le2,\quad 0<s_o\le1.
\]
Write the iterated derivatives as
\[
D^0B=B,\qquad D^{\iota_{\mathrm{der}}+1}B=D(D^{\iota_{\mathrm{der}}}B),
\qquad \iota_{\mathrm{der}}\in\mathbb N,
\]
and define their finite operator-norm bounds by
\[
C_{B,\iota_{\mathrm{der}}}=\sup_{w\in\mathbb R^d}\|D^{\iota_{\mathrm{der}}}B(w)\|,
\qquad \iota_{\mathrm{der}}\in\mathbb N.
\]
Their finiteness follows from continuity and compact support. Set
\[
A_B=\max\{C_{B,p_*},C_{B,p_*+1}\},
\qquad
M_B=2A_B.
\]
The derivative bound and the mean-value inequality give, for every \(w,w'\in\mathbb R^d\),
\[
\|D^{p_*}B(w)-D^{p_*}B(w')\|
\le A_B\|w-w'\|_2,
\qquad
\|D^{p_*}B(w)-D^{p_*}B(w')\|\le2A_B.
\]
If \(\|w-w'\|_2\le1\), use
\(\|w-w'\|_2\le\|w-w'\|_2^{s_o}\); if
\(\|w-w'\|_2>1\), use \(1\le\|w-w'\|_2^{s_o}\).
Thus both cases yield
\[
\|D^{p_*}B(w)-D^{p_*}B(w')\|
\le M_B\|w-w'\|_2^{s_o}.
\]

Translation and dilation satisfy
\[
D^{\iota_{\mathrm{der}}}B_h(x)=h^{-\iota_{\mathrm{der}}}D^{\iota_{\mathrm{der}}}B\bigl((x-x_0)/h\bigr).
\]
Evaluating these multilinear maps in coordinate directions gives the bounds
\[
|\partial^\kappa B_h(x)|
\le d^{\iota_{\mathrm{der}}} C_{B,\iota_{\mathrm{der}}}h^{-\iota_{\mathrm{der}}}
\qquad(|\kappa|=\iota_{\mathrm{der}}\le p_*),
\]
and
\[
|\partial^\kappa B_h(x)-\partial^\kappa B_h(x')|
\le d^{p_*}M_Bh^{-\gamma}\|x-x'\|_2^{s_o}
\qquad(|\kappa|=p_*).
\]
Here the dimension factors are valid because every coordinate direction has norm at most \(d\). Define
\[
S_B=\sum_{\iota_{\mathrm{der}}=0}^{p_*}d^{\iota_{\mathrm{der}}}C_{B,\iota_{\mathrm{der}}},
\qquad
D_B=S_B+d^{p_*}M_B+1.
\]
For \(0<h\le1\) and \(\iota_{\mathrm{der}}\le p_*<\gamma\), one has
\(h^{-\iota_{\mathrm{der}}}\le h^{-\gamma}\). Applying the norm definition in
\cref{obj:def:holder-norm} to the preceding derivative and modulus bounds proves
\[
D_B>0,
\qquad
\|B_h\|_{H^\gamma}\le D_Bh^{-\gamma}
\qquad(0<h\le1).
\]
This also covers the integer smoothness endpoints, where \(s_o=1\).
% lean: CausalSmith.Stat.TwosamplePointcateAnnotationFrontier.oracle_macro_holder

\item Choose the public amplitude constant
\[
\eta_o=\frac{1}{128(D_B+1)}.
\]
It satisfies
\[
\eta_o>0,\qquad
\eta_o(D_B+1)=\frac1{128},\qquad
\eta_o\le\frac18,\qquad
\eta_o\le\frac1{16},\qquad
\eta_oD_B\le L.
\]
For each \(n\ge2\), define
\[
h_n=n^{-1/(2\gamma+d)},
\qquad
b_n=\eta_o r_o(n).
\]
Since \(2\gamma+d>0\) and \(n\ge2\), the power identities give
\[
0<h_n\le1,
\qquad
h_n^\gamma=r_o(n),
\qquad
n\,r_o(n)^2h_n^d=1,
\qquad
0<b_n\le\frac18.
\]
Consequently,
\[
\|b_nB_{h_n}\|_{H^\gamma}
\le b_nD_Bh_n^{-\gamma}
=\eta_oD_B
\le L.
\]

For \(\theta\in\{-1,+1\}\), define the primitive law \(P_{o,\theta}\) by
\[
X\sim\operatorname{Unif}([0,1]^d),
\]
\[
\mathcal L_{P_{o,\theta}}\bigl(A,Y(0),Y(1)\mid X=x\bigr)
=
\operatorname{Bernoulli}(1/2)
\otimes
\operatorname{Bernoulli}(1/2)
\otimes
\operatorname{Bernoulli}\bigl(1/2+\theta b_nB_{h_n}(x)\bigr),
\qquad x\in[0,1]^d.
\]
These conditional margins are Borel. Conditional independence gives exchangeability, and their bounds give
\[
e(x)=\mu_0(x)=\frac12,
\qquad
\frac18\le\frac12+\theta b_nB_{h_n}(x)\le\frac78.
\]
The constant propensity and control functions have Hölder norms at most \(1/2\), while multiplication by either sign preserves the preceding contrast-norm bound. Thus the construction satisfies every condition of \cref{obj:def:primitive-class}.

The designated versions agree with these admissible margins on the cube. Indeed, conditional-margin uniqueness gives almost-sure equality, and continuity together with the full support of the uniform design promotes that equality to pointwise equality. Hence
\[
P_{o,\theta}\in\mathcal M_{1/2},
\qquad
e_{P_{o,\theta}}(x)=\mu_{0,P_{o,\theta}}(x)=\frac12,
\qquad
\tau_{P_{o,\theta}}(x)=\theta b_nB_{h_n}(x)
\quad(x\in[0,1]^d),
\]
so that
\[
\tau_{P_{o,\theta}}(x_0)=\theta b_n.
\]
% lean: CausalSmith.Stat.TwosamplePointcateAnnotationFrontier.oracle_bump_testing_family

\item Define the observed-record and full-experiment laws by
\[
P_{o,\theta}^{\mathrm{obs}}
=\mathcal L_{P_{o,\theta}}(X,A,Y),
\qquad
M_{o,\theta}=\mathsf E_{n,m}(P_{o,\theta}).
\]
Use the common probability reference measure
\[
\nu_o
=
\operatorname{Unif}([0,1]^d)
\otimes\operatorname{Bernoulli}(1/2)
\otimes\operatorname{Bernoulli}(1/2)
\]
on the covariate, treatment, and observed-response coordinates. Relative to this reference, the observed-record density is
\[
\ell_{o,\theta}(x,a_{\mathrm{obs}},y_{\mathrm{obs}})
=
\begin{cases}
1,&a_{\mathrm{obs}}=0,\\
1+2(2y_{\mathrm{obs}}-1)\theta b_nB_{h_n}(x),
&a_{\mathrm{obs}}=1,
\end{cases}
\]
for \(x\in[0,1]^d\) and \(a_{\mathrm{obs}},y_{\mathrm{obs}}\in\{0,1\}\). The amplitude bound implies
\[
\frac12\le\ell_{o,\theta}\le2.
\]
For any two density values at least \(1/2\), the squared sum of their square roots is at least one. The difference-of-squares identity therefore gives, pointwise,
\[
\bigl(\sqrt{\ell_{o,-1}}-\sqrt{\ell_{o,+1}}\bigr)^2
\le
(\ell_{o,-1}-\ell_{o,+1})^2
\le16b_n^2B_{h_n}(x)^2.
\]
The Hellinger normalization in \cref{obj:def:hellinger} can be evaluated under \(\nu_o\): the densities relative to the sum of the two laws are
\(\ell_{o,\theta}/(\ell_{o,-1}+\ell_{o,+1})\), and the common denominator cancels. Thus
\[
\mathsf h^2(P_{o,-1}^{\mathrm{obs}},P_{o,+1}^{\mathrm{obs}})
=
\int
\bigl(\sqrt{\ell_{o,-1}}-\sqrt{\ell_{o,+1}}\bigr)^2\,d\nu_o.
\]
Since \(B_{h_n}^2\le\mathbf1_{C_{h_n}}\), the containment
\(C_{h_n}\subseteq[0,1]^d\) and the product volume of this cube give
\[
\int_{[0,1]^d}B_{h_n}(x)^2\,dx
\le\operatorname{vol}(C_{h_n})
=h_n^d.
\]
Combining the preceding displays proves
\[
\mathsf h^2(P_{o,-1}^{\mathrm{obs}},P_{o,+1}^{\mathrm{obs}})
\le16b_n^2h_n^d.
\]

The auxiliary covariate-treatment marginal is identical under the two laws:
\[
P_{o,-1,XA}=P_{o,+1,XA}
=
\operatorname{Unif}([0,1]^d)\otimes\operatorname{Bernoulli}(1/2).
\]
The product bound and common-factor identity in
\cref{obj:lem:hellinger-block-calculus} therefore yield
\[
\mathsf h^2(M_{o,-1},M_{o,+1})
\le
n\,\mathsf h^2(P_{o,-1}^{\mathrm{obs}},P_{o,+1}^{\mathrm{obs}})
\le16n b_n^2h_n^d
=16\eta_o^2
\le\frac1{16}.
\]
This applies to every \(m\ge0\); at \(m=0\), the auxiliary product is the probability law on the empty record block. The independent uniform randomizer is another common probability factor. Finally, symmetry of the squared square-root difference gives the orientation used below:
\[
\mathsf h^2(M_{o,+1},M_{o,-1})
=
\mathsf h^2(M_{o,-1},M_{o,+1})
\le\frac1{16}.
\]
% lean: CausalSmith.Stat.TwosamplePointcateAnnotationFrontier.oracle_bump_information_at_rate

\item The target is a well-defined functional of these full-experiment laws. To see this, equality of two such laws implies equality of their first labeled-record marginals, since \(n>0\). Under either constructed law, the observed conditional probabilities identify the arm means through
\[
\mu_{1,P_{o,\theta}}(x)
=
2P_{o,\theta}(A=1,Y=1\mid X=x),
\qquad
\mu_{0,P_{o,\theta}}(x)
=
2P_{o,\theta}(A=0,Y=1\mid X=x)
\]
almost everywhere under the uniform design. Thus equality of the observed laws identifies all three conditional margins. Their conditional product construction then identifies the primitive laws, and uniqueness of the canonical continuous contrast gives
\[
M_{o,\theta}=M_{o,\theta'}
\quad\Longrightarrow\quad
\tau_{P_{o,\theta}}(x_0)=\tau_{P_{o,\theta'}}(x_0).
\]
In particular, \(M_{o,+1}\) and \(M_{o,-1}\) are distinct because \(b_n>0\). Define a functional on all measures on the full sample space by
\[
\Psi_o(\mathbb Q_o)=
\begin{cases}
b_n,&\mathbb Q_o=M_{o,+1},\\
-b_n,&\mathbb Q_o=M_{o,-1},\\
0,&\mathbb Q_o\notin\{M_{o,+1},M_{o,-1}\}.
\end{cases}
\]
It satisfies
\[
\Psi_o(M_{o,\theta})=\tau_{P_{o,\theta}}(x_0).
\]

The supplied propensity is the same function for both hypotheses, including its extension outside the cube:
\[
e_o(x)=
\begin{cases}
1/2,&x\in[0,1]^d,\\
0,&x\notin[0,1]^d.
\end{cases}
\]
For any admissible supplied-propensity decision, its section at \(e_o\) is measurable. Apply \cref{obj:lem:published-fuzzy-testing}, the absolute-loss bound of
\citep[Section 3, Lemma 1]{KennedyBalakrishnanRobinsWasserman2024},
with singleton priors, one observation consisting of the entire record block and randomizer, and the choices
\[
\omega_o=\operatorname{Dirac}(*),
\qquad
\mathbb B_*^{\,o}=M_{o,+1},
\qquad
\mathbb C_*^{\,o}=M_{o,-1},
\qquad
\mathcal P_o^{\mathrm{block}}=\{M_{o,+1},M_{o,-1}\},
\]
\[
v_{\mathrm{test}}=1,
\qquad
s_{\mathrm{test}}=2b_n,
\qquad
u_{\mathrm{test}}=\frac1{16}.
\]
The separation is \(2b_n>0\), and the preceding Hellinger bound supplies the distance hypothesis. Consequently,
\[
\sup_{\theta\in\{-1,+1\}}
\mathcal R^e_{n,m}(T^e,P_{o,\theta})
\ge
\frac{b_n}{2}
\left(1-\frac{\sqrt{63}}{32}\right).
\]
Since \(63\le64\),
\[
1-\frac{\sqrt{63}}{32}\ge\frac34,
\qquad
\frac{b_n}{2}
\left(1-\frac{\sqrt{63}}{32}\right)\ge\frac38b_n.
\]
The two primitive laws belong to \(\mathcal M_{1/2}\), so taking the infimum over decisions proves
\[
c=\frac38\eta_o>0,
\qquad
c\,r_o(n)\le R^e_{1/2}(n,m)
\qquad(n\ge2,\ m\ge0).
\]
% lean: CausalSmith.Stat.TwosamplePointcateAnnotationFrontier.oracle_subfamily_lower_rate

\item We next establish the upper bound. For every \(h>0\), define the localization measure and weight by
\[
d\nu_h(x)=h^{-d}\mathbf1_{C_h}(x)\,dx,
\qquad
w_h(x)=h^{-d}\mathbf1_{C_h}(x).
\]
The measure \(\nu_h\) is a probability measure because
\(\operatorname{vol}(C_h)=h^d\). For \(0<h\le1\), its support lies in the design cube.

To make the coarse basis explicit, define
\[
\mathcal I_2=
\{\kappa\in\mathbb N^d:|\kappa|\le2\},
\qquad
q=|\mathcal I_2|,
\]
and the one-dimensional polynomials
\[
\mathsf P_0(\upsilon)=1,\qquad
\mathsf P_1(\upsilon)=\sqrt{12}\,\upsilon,\qquad
\mathsf P_2(\upsilon)=\sqrt5\,(6\upsilon^2-1/2),
\qquad \upsilon\in\mathbb R.
\]
At the current localization width \(h\), the basis and evaluation vector are
\[
r_\kappa(x)=
\prod_{\iota_{\mathrm{cov}}=1}^d
\mathsf P_{\kappa_{\iota_{\mathrm{cov}}}}\bigl((x_{\iota_{\mathrm{cov}}}-1/2)/h\bigr),
\qquad
r(x)=(r_\kappa(x))_{\kappa\in\mathcal I_2},
\qquad
r_0=r(x_0),
\]
for every \(x\in\mathbb R^d\). Define the evaluation kernel and its dimension constant by
\[
\mathcal K_h(x)=r_0^\top r(x),
\qquad
K_{\mathrm{ev}}=q(4^d)^2.
\]
The constant coordinate belongs to \(\mathcal I_2\), so \(q\ge1\) and \(K_{\mathrm{ev}}>0\).

Direct integration of the three one-dimensional formulas gives
\[
\int_{-1/2}^{1/2}
\mathsf P_{j_{\mathrm{pol}}}(\upsilon)
\mathsf P_{j'_{\mathrm{pol}}}(\upsilon)\,d\upsilon
=
\mathbf1_{\{j_{\mathrm{pol}}=j'_{\mathrm{pol}}\}}
\qquad
(j_{\mathrm{pol}},j'_{\mathrm{pol}}\in\{0,1,2\}).
\]
Scaling each coordinate and multiplying these identities yields
\[
\int r_\kappa(x)r_{\kappa'}(x)\,d\nu_h(x)
=
\mathbf1_{\{\kappa=\kappa'\}}.
\]
Expanding the finite sums consequently gives polynomial reproduction and kernel energy:
\[
\int\mathcal K_h(x)\,r(x)^\top v\,d\nu_h(x)
=r_0^\top v
\qquad(v\in\mathbb R^q),
\]
\[
\int\mathcal K_h(x)^2\,d\nu_h(x)=\|r_0\|_2^2.
\]
The displayed polynomial formulas also give
\(|\mathsf P_{j_{\mathrm{pol}}}(\upsilon)|\le4\) on \([-1/2,1/2]\). Thus
\[
|r_\kappa(x)|\le4^d\quad(x\in C_h),
\qquad
|\mathcal K_h(x)|\le K_{\mathrm{ev}}\quad(x\in C_h),
\qquad
\|r_0\|_2^2\le K_{\mathrm{ev}}.
\]
All these identities hold for every \(h>0\).
% lean: CausalSmith.Stat.TwosamplePointcateAnnotationFrontier.oracle_polynomial_reproduction

\item Fix \(P\in\mathcal M\). Define the supplied-propensity response, its localized record contribution, the raw average, and its mean by
\[
Z_P(x,a_{\mathrm{obs}},y_{\mathrm{obs}})
=
\frac{a_{\mathrm{obs}}y_{\mathrm{obs}}}{e_P(x)}
-
\frac{(1-a_{\mathrm{obs}})y_{\mathrm{obs}}}{1-e_P(x)},
\]
\[
F_{P,h}(x,a_{\mathrm{obs}},y_{\mathrm{obs}})
=
w_h(x)\mathcal K_h(x)
Z_P(x,a_{\mathrm{obs}},y_{\mathrm{obs}}),
\]
\[
\overline T_{P,h}(\mathcal D_{n,m},U)
=
\frac1n\sum_{i=1}^nF_{P,h}(X_i,A_i,Y_i),
\qquad
\mu_{P,h}
=
\int\overline T_{P,h}\,d\mathsf E_{n,m}(P).
\]
These definitions apply to all \(x\in\mathbb R^d\), binary record coordinates, and every dataset and randomizer, using the quotient convention stated above.

On the design cube, overlap and the binary record coordinates give
\[
|Z_P(x,a_{\mathrm{obs}},y_{\mathrm{obs}})|\le\varepsilon^{-1}.
\]
Outside the cube, the zero extension of the propensity gives
\[
Z_P(x,a_{\mathrm{obs}},y_{\mathrm{obs}})
=-(1-a_{\mathrm{obs}})y_{\mathrm{obs}},
\qquad
|Z_P(x,a_{\mathrm{obs}},y_{\mathrm{obs}})|\le1\le\varepsilon^{-1}.
\]
Hence the response bound is global, and
\[
|F_{P,h}(x,a_{\mathrm{obs}},y_{\mathrm{obs}})|
\le\varepsilon^{-1}h^{-d}K_{\mathrm{ev}}.
\]
Borel measurability and this bound supply all integrability and finite second moments used below.

For \(0<h\le1\), exchangeability and the conditional-margin identities in
\cref{obj:def:primitive-class}, followed by cancellation of the overlap denominators, give
\[
\mathbb E_P\left[
w_h(X)\mathcal K_h(X)\frac{AY}{e_P(X)}
\right]
=
\int_{[0,1]^d}w_h(x)\mathcal K_h(x)\mu_{1,P}(x)\,dx,
\]
\[
\mathbb E_P\left[
w_h(X)\mathcal K_h(X)\frac{(1-A)Y}{1-e_P(X)}
\right]
=
\int_{[0,1]^d}w_h(x)\mathcal K_h(x)\mu_{0,P}(x)\,dx.
\]
Subtracting and using the canonical contrast identity gives
\[
\mathbb E_P[F_{P,h}(X,A,Y)]
=
\int\mathcal K_h(x)\tau_P(x)\,d\nu_h(x).
\]
Every labeled coordinate has this expectation, while the auxiliary block and randomizer integrate out as probability factors. Since \(n>0\),
\[
\mu_{P,h}
=
\int\mathcal K_h(x)\tau_P(x)\,d\nu_h(x).
\]
% lean: CausalSmith.Stat.TwosamplePointcateAnnotationFrontier.oracle_average_mean

\item We now bound this mean's bias over the full range \(0<h\le1\). Define the Taylor polynomial
\[
p(x)=
\sum_{|\kappa|\le p_*}
\frac{\partial^\kappa\tau_P(x_0)}{\kappa!}
(x-x_0)^\kappa,
\qquad
\kappa!=\prod_{\iota_{\mathrm{cov}}=1}^d\kappa_{\iota_{\mathrm{cov}}}!.
\]
Because \(p_*\le2\), the coarse polynomial basis spans \(p\). Thus there exists a coefficient vector \(\vartheta_{\mathrm{Tay}}\in\mathbb R^q\) with
\[
p_{\vartheta_{\mathrm{Tay}}}(x)
=
r(x)^\top\vartheta_{\mathrm{Tay}}
=
p(x)
\qquad(x\in\mathbb R^d),
\qquad
p_{\vartheta_{\mathrm{Tay}}}(x_0)=\tau_P(x_0).
\]

If \(\gamma=1\), then \(p_*=0\) and \(p(x)=\tau_P(x_0)\). The Hölder norm condition gives the Lipschitz bound, and the geometry of \(C_h\) gives
\[
|\tau_P(x)-p(x)|
\le L\|x-x_0\|_2
\le L\sqrt d\,h
\le L(\sqrt d+1)h
\qquad(x\in C_h).
\]

If \(1<\gamma\le3\), then \(1\le p_*\le2\). For each \(x\in C_h\), define the segment function
\[
g_{P,x}(\sigma_{\mathrm{seg}})=\tau_P\bigl(x_0+\sigma_{\mathrm{seg}}(x-x_0)\bigr),
\qquad 0\le\sigma_{\mathrm{seg}}\le1.
\]
The segment lies in the design cube. Expanding its \(p_*\)-th derivative in coordinate partials and using their \(s_o\)-Hölder modulus gives
\[
\left|g_{P,x}^{(p_*)}(\sigma_{\mathrm{seg}})-g_{P,x}^{(p_*)}(0)\right|
\le
(d+1)^{p_*}L
\|x-x_0\|_2^{p_*+s_o}\sigma_{\mathrm{seg}}^{s_o}.
\]
Indeed, the multinomial expansion contributes at most
\((d+1)^{p_*}\|x-x_0\|_2^{p_*}\), and each top-partial difference is bounded by
\(L\sigma_{\mathrm{seg}}^{s_o}\|x-x_0\|_2^{s_o}\).
The one-dimensional Lagrange remainder at degree \(p_*-1\), followed by subtraction of the degree-\(p_*\) Taylor term, supplies some
\(\xi_{P,x}\in(0,1)\) such that
\[
\tau_P(x)-p(x)
=
\frac{
g_{P,x}^{(p_*)}(\xi_{P,x})
-g_{P,x}^{(p_*)}(0)
}{p_*!}.
\]
Since \(\xi_{P,x}^{s_o}\le1\) and \(p_*!\ge1\),
\[
|\tau_P(x)-p(x)|
\le(d+1)^{p_*}L\|x-x_0\|_2^\gamma.
\]
Using \(\|x-x_0\|_2\le(d+1)h\) and
\(\gamma=p_*+s_o\le p_*+1\) yields
\[
|\tau_P(x)-p(x)|
\le L(d+1)^{2p_*+1}h^\gamma
\qquad(x\in C_h).
\]

Define the positive public constants
\[
C_{\mathrm{Tay}}=
\begin{cases}
L(\sqrt d+1),&\gamma=1,\\
L(d+1)^{2p_*+1},&1<\gamma\le3,
\end{cases}
\qquad
C_{\mathrm{bias}}=K_{\mathrm{ev}}C_{\mathrm{Tay}}.
\]
Both cases establish
\[
\sup_{x\in C_h}|\tau_P(x)-p_{\vartheta_{\mathrm{Tay}}}(x)|
\le C_{\mathrm{Tay}}h^\gamma.
\]
Polynomial reproduction and the mean identity therefore imply
\[
\begin{aligned}
|\mu_{P,h}-\tau_P(x_0)|
&=
\left|
\int\mathcal K_h(x)
\bigl(\tau_P(x)-p_{\vartheta_{\mathrm{Tay}}}(x)\bigr)
\,d\nu_h(x)
\right|\\
&\le
\int|\mathcal K_h(x)|
\left|\tau_P(x)-p_{\vartheta_{\mathrm{Tay}}}(x)\right|
\,d\nu_h(x)\\
&\le C_{\mathrm{bias}}h^\gamma.
\end{aligned}
\]
At \(h=h_n\), this becomes
\[
|\mu_{P,h_n}-\tau_P(x_0)|
\le C_{\mathrm{bias}}r_o(n).
\]
% lean: CausalSmith.Stat.TwosamplePointcateAnnotationFrontier.oracle_tuned_average_bias

\item The global response bound gives the recordwise squared bound
\[
F_{P,h}(X,A,Y)^2
\le\varepsilon^{-2}\bigl(w_h(X)\mathcal K_h(X)\bigr)^2.
\]
Under the uniform design, localization and the kernel-energy identity give
\[
\begin{aligned}
\mathbb E_P[F_{P,h}(X,A,Y)^2]
&\le
\varepsilon^{-2}\int_{[0,1]^d}
\bigl(w_h(x)\mathcal K_h(x)\bigr)^2\,dx\\
&=
\varepsilon^{-2}h^{-d}\int\mathcal K_h(x)^2\,d\nu_h(x)\\
&=
\varepsilon^{-2}h^{-d}\|r_0\|_2^2.
\end{aligned}
\]
Independence of the \(n\) labeled records yields
\[
\begin{aligned}
\int(\overline T_{P,h}-\mu_{P,h})^2\,d\mathsf E_{n,m}(P)
&=
\frac1n\operatorname{Var}_P(F_{P,h}(X,A,Y))\\
&\le\frac1n\mathbb E_P[F_{P,h}(X,A,Y)^2]\\
&\le\frac{\varepsilon^{-2}\|r_0\|_2^2}{nh^d}.
\end{aligned}
\]
The first inequality uses
\[
\operatorname{Var}_P(F_{P,h})
=
\mathbb E_P[F_{P,h}^2]
-\bigl(\mathbb E_P[F_{P,h}]\bigr)^2
\le\mathbb E_P[F_{P,h}^2].
\]
Cauchy--Schwarz on the probability experiment then gives
\[
\int|\overline T_{P,h}-\mu_{P,h}|\,d\mathsf E_{n,m}(P)
\le
\left(
\int(\overline T_{P,h}-\mu_{P,h})^2\,d\mathsf E_{n,m}(P)
\right)^{1/2}
\le
\sqrt{\frac{\varepsilon^{-2}\|r_0\|_2^2}{nh^d}}.
\]
At \(h=h_n\), use the target-coordinate bound and
\[
(nh_n^d)^{-1/2}=r_o(n)
\]
to obtain
\[
\int|\overline T_{P,h_n}-\mu_{P,h_n}|\,d\mathsf E_{n,m}(P)
\le\sqrt{\varepsilon^{-2}K_{\mathrm{ev}}}\,r_o(n)
\le C_{\mathrm{dev}}r_o(n),
\qquad
C_{\mathrm{dev}}=\sqrt{\varepsilon^{-2}K_{\mathrm{ev}}}+1>0.
\]
The estimate is uniform in \(P\), \(n\ge2\), and \(m\ge0\).
% lean: CausalSmith.Stat.TwosamplePointcateAnnotationFrontier.oracle_tuned_absolute_deviation

\item The probability-valued arm-mean conditions in \cref{obj:def:primitive-class} imply
\[
-1
\le
\mu_{1,P}(x_0)-\mu_{0,P}(x_0)
=
\tau_P(x_0)
\le1.
\]
The displayed smoother equals
\[
T^e(\mathcal D_{n,m},U,e_P)
=
\operatorname{clip}_{[-1,1]}
\bigl(\overline T_{P,h_n}(\mathcal D_{n,m},U)\bigr).
\]
Projection onto this interval contracts distance to any target inside it. Thus, for every dataset and randomizer,
\[
\begin{aligned}
|T^e(\mathcal D_{n,m},U,e_P)-\tau_P(x_0)|
&\le|\overline T_{P,h_n}(\mathcal D_{n,m},U)-\tau_P(x_0)|\\
&\le
|\overline T_{P,h_n}(\mathcal D_{n,m},U)-\mu_{P,h_n}|
+
|\mu_{P,h_n}-\tau_P(x_0)|.
\end{aligned}
\]
Integrating and applying the preceding two bounds gives
\[
\mathcal R^e_{n,m}(T^e,P)
\le
(C_{\mathrm{dev}}+C_{\mathrm{bias}})r_o(n).
\]
Define
\[
C_{\mathrm{up}}=C_{\mathrm{dev}}+C_{\mathrm{bias}}>0.
\]
The bounded measurable raw average and continuous clipping map give measurability of the supplied-propensity section. The clipping formula itself gives the pointwise guarantee
\[
\left|T^e(\mathcal D_{n,m},U,e_P)\right|\le1
\]
for every dataset and randomizer.
% lean: CausalSmith.Stat.TwosamplePointcateAnnotationFrontier.oracle_smoother_upper_rate

\item For completeness, the smoother has a total admissible extension to arbitrary supplied functions. Define
\[
\widetilde T^e(\mathcal D_{n,m},U,e)=
\begin{cases}
T^e(\mathcal D_{n,m},U,e),&e\text{ is Borel},\\
0,&e\text{ is not Borel}.
\end{cases}
\]
For every fixed Borel \(e\), the totalized reciprocal, record-coordinate maps, finite sum, and clipping are Borel; the other section is constant. Since each designated \(e_P\) is Borel, this extension agrees with the displayed smoother at every admitted law. It is therefore a candidate in \cref{obj:def:oracle-risk}, and the uniform bound proves
\[
R^e(n,m)
\le
\sup_{P\in\mathcal M}
\mathcal R^e_{n,m}(\widetilde T^e,P)
\le C_{\mathrm{up}}r_o(n).
\]

Finally, retain the lower constant already constructed and set
\[
C=\max\{c,C_{\mathrm{up}}\}.
\]
These are finite real constants depending on the fixed public parameters, with
\[
0<c\le C.
\]
Because \(r_o(n)\ge0\), enlarging the upper constant preserves both the minimax and individual-decision bounds. Combining the lower bound with subfamily monotonicity gives, for every \(n\ge2\) and \(m\ge0\),
\[
c\,r_o(n)
\le R^e_{1/2}(n,m)
\le R^e(n,m)
\le C\,r_o(n),
\]
and, for every \(P\in\mathcal M\),
\[
\mathcal R^e_{n,m}(T^e,P)\le C\,r_o(n).
\]
Together with the measurability and pointwise clipping bound above, these are all the asserted conclusions.
% lean: CausalSmith.Stat.TwosamplePointcateAnnotationFrontier.supplied_propensity
\end{enumerate}
\end{proof}

\begin{proof}[Proof of \cref{obj:thm:marked-component-certificate}]
Fix the public parameters. All constants constructed below depend only on these parameters. We use the functions and sign arrays prescribed by \cref{obj:def:marked-handle}, extending the displayed function formulas to \(\mathbb R^d\) when taking analytic bounds.

\begin{enumerate}
\item \textit{Regularity and membership.}
Set
\[
D_d:=2^d.
\]
The function \(\chi\) is smooth and flat at zero: each of its derivatives on the positive half-line is a polynomial in \(1/u\) multiplied by \(\exp(-1/u)\), and this expression tends to zero as \(u\downarrow0\). The denominator defining the scalar angular transition in \cref{obj:def:marked-handle} is positive everywhere. Consequently the prescribed sine and cosine pieces join smoothly, including at their support endpoints. The bump \(B\) is also smooth and compactly supported, and
\[
0\le B_h(x)\le1,\qquad B_h(x_0)=1.
\]
Define the finite derivative bounds
\[
K_\omega:=1+\sup_{u\in\mathbb R}|\omega_0'(u)|,
\qquad
K_B:=1+\sup_{w\in\mathbb R^d}\|\nabla B(w)\|_2,
\]
and put
\[
K_\phi:=K_B+dK_\omega,
\qquad
K_F:=2D_dK_\phi.
\]
Translation invariance gives the same derivative bound for every \(\omega_z\). Telescoping a product of factors bounded by one, and then using \(\delta\le h\), yields
\[
|\phi_{\mathbf z}(x)-\phi_{\mathbf z}(x')|
\le
\left(\frac{K_B}{h}+\frac{dK_\omega}{\delta}\right)
\|x-x'\|_2
\le
\frac{K_\phi}{\delta}\|x-x'\|_2.
\]

At each point, a nonzero frame entry has, in each coordinate, one of the two lattice indices adjacent to the scaled coordinate. Thus at most \(D_d\) entries are nonzero. Each entry has absolute value at most one, so
\[
|F_\lambda(x)|\le D_d,
\qquad
|G_\eta(x)|\le D_d.
\]
For a difference between two points, entries outside the union of their active index sets contribute zero. That union has at most \(2D_d\) elements. Therefore
\[
|F_\lambda(x)-F_\lambda(x')|
\le \frac{K_F}{\delta}\|x-x'\|_2,
\qquad
|G_\eta(x)-G_\eta(x')|
\le \frac{K_F}{\delta}\|x-x'\|_2.
\]
For a Hölder order \(s_{\mathrm{ord}}\in(0,1]\), define
\[
C_{F,s_{\mathrm{ord}}}:=2D_d+K_F+1.
\]
When \(\|x-x'\|_2\le\delta\), the Lipschitz bound and
\[
\frac{\|x-x'\|_2}{\delta}
\le
\left(\frac{\|x-x'\|_2}{\delta}\right)^{s_{\mathrm{ord}}}
\]
give the required Hölder difference bound. When \(\|x-x'\|_2>\delta\), use the bound \(2D_d\) instead. Since \(\delta\le1\), the supremum term is covered as well. With the norm of \cref{obj:def:holder-norm}, this proves
\[
\|F_\lambda\|_{H^{s_{\mathrm{ord}}}}
\le C_{F,s_{\mathrm{ord}}}\delta^{-s_{\mathrm{ord}}},
\qquad
\|G_\eta\|_{H^{s_{\mathrm{ord}}}}
\le C_{F,s_{\mathrm{ord}}}\delta^{-s_{\mathrm{ord}}}.
\]

For the macro bump, use its fixed compact profile \(B^2\). For every integer \(j\ge0\), define
\[
D_{B,j}:=
1+\sup_{w\in\mathbb R^d}\|\mathrm D^j(B^2)(w)\|_{\mathrm{op}},
\]
where \(\mathrm D^j\) is the total derivative of order \(j\), with the order-zero norm interpreted as absolute value. These constants are finite by smoothness and compact support. For \(s_{\mathrm{ord}}>0\), define
\[
j_{\mathrm{der}}:=\lceil s_{\mathrm{ord}}\rceil-1,
\qquad
\xi_{\mathrm{der}}:=s_{\mathrm{ord}}-j_{\mathrm{der}}\in(0,1],
\]
\[
M_{B,s_{\mathrm{ord}}}
:=
2\max\{D_{B,j_{\mathrm{der}}},D_{B,j_{\mathrm{der}}+1}\},
\]
\[
C_{B,s_{\mathrm{ord}}}
:=
\sum_{j=0}^{j_{\mathrm{der}}}D_{B,j}d^j
+
M_{B,s_{\mathrm{ord}}}d^{j_{\mathrm{der}}}+1.
\]
The derivative of order \(j_{\mathrm{der}}\) has both a bounded supremum and a Lipschitz bound supplied by the next derivative. Splitting at distance one gives
\[
\|\mathrm D^{j_{\mathrm{der}}}(B^2)(w)
-\mathrm D^{j_{\mathrm{der}}}(B^2)(w')\|_{\mathrm{op}}
\le
M_{B,s_{\mathrm{ord}}}\|w-w'\|_2^{\xi_{\mathrm{der}}}.
\]
Dilation by \(h\) multiplies a derivative of order \(j\) by \(h^{-j}\). Evaluating a derivative on coordinate directions bounds a partial derivative by its operator norm times \(d^j\). Hence, for \(0<h\le1\),
\[
|\partial^\kappa(B_h^2)(x)|
\le D_{B,|\kappa|}d^{|\kappa|}h^{-|\kappa|}
\le C_{B,s_{\mathrm{ord}}}h^{-s_{\mathrm{ord}}}
\qquad (|\kappa|\le j_{\mathrm{der}}),
\]
and
\[
|\partial^\kappa(B_h^2)(x)-\partial^\kappa(B_h^2)(x')|
\le
M_{B,s_{\mathrm{ord}}}d^{j_{\mathrm{der}}}
h^{-s_{\mathrm{ord}}}\|x-x'\|_2^{\xi_{\mathrm{der}}}
\qquad (|\kappa|=j_{\mathrm{der}}).
\]
Thus, in particular for the orders \(\beta\) and \(\gamma\),
\[
\|B_h^2\|_{H^{s_{\mathrm{ord}}}}
\le C_{B,s_{\mathrm{ord}}}h^{-s_{\mathrm{ord}}}.
\]

Define
\[
K_M:=
D_d+C_{F,\alpha}+C_{F,\beta}
+C_{B,\gamma}+C_{B,\beta}+1,
\qquad
\rho:=\min\left\{1,L-1/2,\frac12-\varepsilon\right\},
\qquad
c_M:=\frac{\rho}{32K_M}.
\]
The public domain gives \(0<\rho\le1\). Hence \(0<c_M\le\rho/32\le1/32\), and multiplying any of the five summands preceding the final \(1\) in \(K_M\) by \(c_M\) gives at most \(\rho/32\).

Suppose the amplitude and contrast inequalities hold with \(c_M\). Since \(0<\delta\le h\le1/2\),
\[
a\le c_M,\qquad b\le c_M,\qquad ab\le c_M,
\]
\[
a\delta^{-\alpha}\le c_M,\qquad
b\delta^{-\beta}\le c_M,\qquad
ab h^{-\gamma}\le c_M.
\]
Also \(\beta\le\gamma\) and \(h\le1\), so
\[
ab\le c_Mh^\gamma\le c_Mh^\beta,
\qquad
ab h^{-\beta}\le c_M.
\]
The prescribed probabilities therefore satisfy
\[
|e_\lambda(x)-1/2|
\le aD_d\le\frac{\rho}{32}\le\frac14,
\qquad
aD_d\le\frac{\rho}{32}\le\frac12-\varepsilon,
\]
\[
|\mu_{0,\theta,\eta}(x)-1/2|,
\ |\mu_{1,\theta,\eta}(x)-1/2|
\le bD_d+ab
\le\frac{\rho}{32}+c_M\le\frac1{16}\le\frac38.
\]
These bounds imply the overlap and both probability-valued arm-mean conditions.

Subadditivity and homogeneity of the Hölder norm give
\[
\|aF_\lambda\|_{H^\alpha}\le C_{F,\alpha}c_M,
\qquad
\|bG_\eta\|_{H^\beta}\le C_{F,\beta}c_M,
\]
\[
\|\tau_\theta\|_{H^\gamma}
\le2C_{B,\gamma}c_M,
\qquad
\|\theta abB_h^2\|_{H^\beta}
\le C_{B,\beta}c_M.
\]
Consequently
\[
\|e_\lambda\|_{H^\alpha}
\le\frac12+\frac{\rho}{32}\le L,
\]
\[
\|\mu_{0,\theta,\eta}\|_{H^\beta}
\le\frac12+\frac{\rho}{16}\le L,
\qquad
\|\tau_\theta\|_{H^\gamma}\le\frac{\rho}{16}\le L.
\]
All the prescribed functions are Borel. The independent conditional Bernoulli construction gives uniform covariates and conditional exchangeability, and its conditional margins are exactly the prescribed functions. Together with the preceding bounds, these facts establish every condition in \cref{obj:def:primitive-class}.

To identify the canonical contrast, the prescribed control mean is continuous, and the prescribed treated mean is continuous because it is the sum of the control mean and \(\tau_\theta\). Their admissible designated versions agree with them almost surely by uniqueness of conditional margins. Two continuous versions that agree almost everywhere under the uniform design agree everywhere on the cube: a nonzero continuous difference would persist on a relative neighborhood of positive uniform measure, including at a boundary point. Thus
\[
P_{\theta,\lambda,\eta}\in\mathcal M,
\qquad
\tau_{P_{\theta,\lambda,\eta}}(x)
=
\mu_{1,\theta,\eta}(x)-\mu_{0,\theta,\eta}(x)
=
\tau_\theta(x)
\quad(x\in[0,1]^d).
\]
% lean: CausalSmith.Stat.TwosamplePointcateAnnotationFrontier.marked_membership

\item \textit{Normalization and singleton cancellation.}
Define expectation over the finite sign arrays by
\[
\mathbb E_\theta Q
:=
\sum_{\lambda,\eta}
w_\theta(\lambda,\eta)Q(\lambda,\eta).
\]
At each index, direct summation over the four sign pairs gives
\[
\frac{1+\theta l s_\eta/2}{4}\ge0
\quad(l,s_\eta\in\{-1,+1\}),
\qquad
\sum_{l,s_\eta\in\{-1,+1\}}\frac{1+\theta l s_\eta/2}{4}=1.
\]
Multiplying these identities over the finite index set proves
\[
w_\theta(\lambda,\eta)\ge0,
\qquad
\sum_{\lambda,\eta}w_\theta(\lambda,\eta)=1.
\]
The same finite-product summation gives
\[
\mathbb E_\theta\lambda_{\mathbf z}
=
\mathbb E_\theta\eta_{\mathbf z}=0,
\qquad
\mathbb E_\theta[\lambda_{\mathbf z}\eta_{\mathbf w}]
=
\begin{cases}
\theta/2,&\mathbf z=\mathbf w,\\
0,&\mathbf z\ne\mathbf w.
\end{cases}
\]
Indeed, matching indices use the scalar sign-pair correlation, while distinct indices include a zero-mean scalar factor.

The sine-cosine identity on each lattice interval gives a sum of squares equal to one. Tensoring it and multiplying by \(B_h^2\) yields
\[
\sum_{\mathbf z\in\mathcal I}\phi_{\mathbf z}(x)^2=B_h(x)^2.
\]
For \(x\in C_h\), all nonzero tensor entries lie in the prescribed finite index set; for \(x\notin C_h\), both sides vanish. Expanding the finite signed sums therefore gives
\[
\mathbb E_\theta F_\lambda(x)
=
\mathbb E_\theta G_\eta(x)=0,
\qquad
\mathbb E_\theta[F_\lambda(x)G_\eta(x)]
=\frac{\theta}{2}B_h(x)^2.
\]

Set
\[
c_S:=\frac1{64D_d}.
\]
If the amplitude inequalities hold with \(c_S\), then
\[
a,b\le c_S,\qquad
aD_d,bD_d\le\frac1{64},
\qquad
ab\le\frac1{4096}.
\]
Thus the raw propensity and both raw arm probabilities lie in \([0,1]\), at every covariate, so their Bernoulli atom formulas apply directly.

For a labeled record, define the treatment and observed-response spins by
\[
s^A:=2A-1,\qquad
Y:=AY(1)+(1-A)Y(0),\qquad
s^Y:=2Y-1.
\]
Relative to the uniform probability measure on the four spin pairs, its likelihood is
\[
L_{\theta,\lambda,\eta}(x;s^A,s^Y)
:=
(1+2s^AaF_\lambda(x))
(1+2s^YbG_\eta(x)+s^As^Y\tau_\theta(x)).
\]
This follows by multiplying the two Bernoulli atom probabilities, each divided by its fair atom probability. Expanding the product gives
\[
\begin{aligned}
\mathbb E_\theta L_{\theta,\lambda,\eta}(x;s^A,s^Y)
={}&1+s^As^Y\tau_\theta(x)\\
&+\bigl(2s^Aa+2(s^A)^2s^Ya\tau_\theta(x)\bigr)
  \mathbb E_\theta F_\lambda(x)\\
&+2s^Yb\,\mathbb E_\theta G_\eta(x)
+4s^As^Yab\,\mathbb E_\theta[F_\lambda(x)G_\eta(x)]\\
={}&1+s^As^Y
\bigl(-2\theta abB_h(x)^2+2\theta abB_h(x)^2\bigr)
=1.
\end{aligned}
\]
Hence every atom of the singleton labeled mixture has probability \(1/4\). This proves its uniform conditional law for every hypothesis and every cube covariate.
% lean: CausalSmith.Stat.TwosamplePointcateAnnotationFrontier.singleton_cancellation

\item \textit{Equality on uninformative components.}
For integers \(n\ge2\) and \(m\ge0\), write
\[
N:=n+m,\qquad
\mathbf x:=(x_1,\ldots,x_N)\in([0,1]^d)^N.
\]
The first \(n\) indices are labeled. A spin assignment has domain
\[
\mathbf s\in
\bigl(\{-1,+1\}^2\bigr)^n
\times\{-1,+1\}^m.
\]
Define the record likelihoods by
\[
R_{\theta,\lambda,\eta,i}(\mathbf x,\mathbf s)
:=
\begin{cases}
L_{\theta,\lambda,\eta}(x_i;s_i^A,s_i^Y),&i\le n,\\
1+2s_i^AaF_\lambda(x_i),&i>n,
\end{cases}
\]
and, for every subset \(\mathcal V\subseteq\{1,\ldots,N\}\), define
\[
g_{\theta,\mathcal V}(\mathbf x,\mathbf s)
:=
\mathbb E_\theta
\prod_{i\in\mathcal V}
R_{\theta,\lambda,\eta,i}(\mathbf x,\mathbf s).
\]
Empty products equal one, so \(g_{\theta,\varnothing}=1\).

The sign flip
\[
\mathcal F_\eta(\lambda,\eta):=(\lambda,-\eta)
\]
is a bijection, and its weights satisfy
\[
w_{+1}(\lambda,\eta)=w_{-1}(\lambda,-\eta).
\]
Every auxiliary likelihood is unchanged by this flip, because it depends on \(\lambda\) alone. Reindexing the finite sum consequently gives
\[
g_{+1,\mathcal V}(\mathbf x,\mathbf s)
=
g_{-1,\mathcal V}(\mathbf x,\mathbf s)
\qquad
\text{if }\mathcal V\subseteq\{n+1,\ldots,N\}.
\]
If \(\mathcal V=\{i\}\) with \(i\le n\), singleton cancellation instead gives
\[
g_{+1,\{i\}}(\mathbf x,\mathbf s)
=
g_{-1,\{i\}}(\mathbf x,\mathbf s)=1.
\]
Thus we may take
\[
c_U:=c_S.
\]
Under its amplitude conditions, every auxiliary-only component and every singleton labeled component has equal densities under the two hypotheses. Equality for all spin assignments gives equality of their full conditional laws.
% lean: CausalSmith.Stat.TwosamplePointcateAnnotationFrontier.uninformative_component_equality

\item \textit{The component density difference.}
Define
\[
c_\Delta:=\frac1{64D_d},
\qquad
C_\Delta:=8D_d^2+6.
\]
Assume the amplitude inequalities with \(c_\Delta\). Fix \(\mathbf x,\mathbf s\) and a subset \(\mathcal V\), and define its size by
\[
p_{\mathrm{cmp}}:=|\mathcal V|.
\]
For every record, introduce the contrast-deleted factors
\[
f_{\lambda,i}:=1+2s_i^AaF_\lambda(x_i),
\qquad
g_{\eta,i}:=
\begin{cases}
1+2s_i^YbG_\eta(x_i),&i\le n,\\
1,&i>n.
\end{cases}
\]
The bounds from the previous steps give
\[
|f_{\lambda,i}-1|\le2aD_d,\qquad
|g_{\eta,i}-1|\le2bD_d,
\qquad
|f_{\lambda,i}|,|g_{\eta,i}|\le\frac98.
\]
For labeled records,
\[
R_{\theta,\lambda,\eta,i}
=f_{\lambda,i}
\bigl(g_{\eta,i}+s_i^As_i^Y\tau_\theta(x_i)\bigr),
\qquad
|\tau_\theta(x_i)|\le2ab.
\]
For auxiliary records, \(R_{\theta,\lambda,\eta,i}=f_{\lambda,i}\). Since \(ab\le1/4096\), both cases satisfy
\[
|R_{\theta,\lambda,\eta,i}|
\le\frac32,
\qquad
|R_{\theta,\lambda,\eta,i}-f_{\lambda,i}g_{\eta,i}|
\le3ab,
\qquad
|f_{\lambda,i}g_{\eta,i}|\le\frac32.
\]

Define
\[
\mathcal A_{\mathcal V}(\lambda)
:=\prod_{i\in\mathcal V}f_{\lambda,i},
\qquad
\mathcal G_{\mathcal V}(\eta)
:=\prod_{i\in\mathcal V}g_{\eta,i},
\]
\[
U_{\mathcal V}
:=p_{\mathrm{cmp}}(2aD_d)(9/8)^{p_{\mathrm{cmp}}},
\qquad
W_{\mathcal V}
:=p_{\mathrm{cmp}}(2bD_d)(9/8)^{p_{\mathrm{cmp}}},
\]
\[
D_{\mathcal V}
:=p_{\mathrm{cmp}}(3ab)(3/2)^{p_{\mathrm{cmp}}}.
\]
Telescoping a product, using an upper bound \(D_{\mathrm{prod}}\ge1\) on both sets of factors, gives
\[
\left|
\prod_{i\in\mathcal V}f_{\mathrm{tel},i}
-\prod_{i\in\mathcal V}g_{\mathrm{tel},i}
\right|
\le
D_{\mathrm{prod}}^{|\mathcal V|}
\sum_{i\in\mathcal V}
|f_{\mathrm{tel},i}-g_{\mathrm{tel},i}|.
\]
Here the telescoping sum replaces one factor at a time; each remaining product is bounded by \(D_{\mathrm{prod}}^{|\mathcal V|-1}\), which is at most the displayed power. For the empty set both sides are zero. Applying this inequality gives
\[
|\mathcal A_{\mathcal V}-1|\le U_{\mathcal V},
\qquad
|\mathcal G_{\mathcal V}-1|\le W_{\mathcal V},
\]
\[
\left|
\prod_{i\in\mathcal V}R_{\theta,\lambda,\eta,i}
-\mathcal A_{\mathcal V}\mathcal G_{\mathcal V}
\right|
\le D_{\mathcal V}.
\]

Flipping \(\eta\) preserves every function of \(\lambda\) and changes the positive weights to the negative weights. Flipping \(\lambda\) similarly preserves every function of \(\eta\). Consequently
\[
\mathbb E_{+1}\mathcal A_{\mathcal V}
=\mathbb E_{-1}\mathcal A_{\mathcal V},
\qquad
\mathbb E_{+1}\mathcal G_{\mathcal V}
=\mathbb E_{-1}\mathcal G_{\mathcal V}.
\]
Expanding the centered product and using normalization therefore yields
\[
\begin{aligned}
&\mathbb E_{+1}[\mathcal A_{\mathcal V}\mathcal G_{\mathcal V}]
-\mathbb E_{-1}[\mathcal A_{\mathcal V}\mathcal G_{\mathcal V}]\\
&\qquad=
\mathbb E_{+1}[(\mathcal A_{\mathcal V}-1)(\mathcal G_{\mathcal V}-1)]
-\mathbb E_{-1}[(\mathcal A_{\mathcal V}-1)(\mathcal G_{\mathcal V}-1)].
\end{aligned}
\]
Each expectation on the right has absolute value at most
\(U_{\mathcal V}W_{\mathcal V}\). Adding the two contrast-deletion errors gives
\[
|g_{+1,\mathcal V}-g_{-1,\mathcal V}|
\le D_{\mathcal V}+2U_{\mathcal V}W_{\mathcal V}+D_{\mathcal V}.
\]
Finally,
\[
(9/8)^{p_{\mathrm{cmp}}}(9/8)^{p_{\mathrm{cmp}}}
\le(3/2)^{p_{\mathrm{cmp}}},
\qquad
p_{\mathrm{cmp}}\le p_{\mathrm{cmp}}^2
\quad(p_{\mathrm{cmp}}\in\mathbb N).
\]
Substitution proves
\[
|g_{+1,\mathcal V}(\mathbf x,\mathbf s)
-g_{-1,\mathcal V}(\mathbf x,\mathbf s)|
\le
C_\Delta(3/2)^{p_{\mathrm{cmp}}}
p_{\mathrm{cmp}}^2ab.
\]
For \(p_{\mathrm{cmp}}=0\), both densities equal one and the bound is zero.
% lean: CausalSmith.Stat.TwosamplePointcateAnnotationFrontier.component_information_bound

\item \textit{A positive floor and component Hellinger distance.}
Define
\[
c_H:=\min\left\{c_\Delta,\frac1{64D_d}\right\},
\qquad
K_{\mathrm{info}}:=C_\Delta^2.
\]
Under the amplitude inequalities with \(c_H\),
\[
aD_d\le\frac1{16},
\qquad
bD_d+ab\le\frac1{16}.
\]
Thus the treatment factor of a labeled likelihood is at least \(7/8\), and its outcome factor is also at least \(7/8\). Their product is at least \(1/2\). An auxiliary likelihood is at least \(7/8\), hence at least \(1/2\). Averaging the resulting product floor with nonnegative normalized weights gives
\[
g_{\theta,\mathcal V}(\mathbf x,\mathbf s)
\ge(1/2)^{p_{\mathrm{cmp}}}.
\]

Define the common probability reference on the record spins by
\[
\mu_{\mathrm{spin}}
:=
\bigotimes_{i=1}^{n}\operatorname{Unif}(\{-1,+1\}^2)
\otimes
\bigotimes_{i=n+1}^{N}\operatorname{Unif}(\{-1,+1\}),
\]
and define
\[
H_{\mathcal V}^2(\mathbf x)
:=
\int
\left(
\sqrt{g_{+1,\mathcal V}(\mathbf x,\mathbf s)}
-\sqrt{g_{-1,\mathcal V}(\mathbf x,\mathbf s)}
\right)^2
\,d\mu_{\mathrm{spin}}(\mathbf s).
\]
For nonnegative numbers bounded below by \(t_{\mathrm{floor}}>0\), the identity
\[
(f_{\mathrm{den}}-g_{\mathrm{den}})^2
=
(\sqrt{f_{\mathrm{den}}}-\sqrt{g_{\mathrm{den}}})^2
(\sqrt{f_{\mathrm{den}}}+\sqrt{g_{\mathrm{den}}})^2
\]
and the denominator bound
\[
(\sqrt{f_{\mathrm{den}}}+\sqrt{g_{\mathrm{den}}})^2
\ge4t_{\mathrm{floor}}
\]
give
\[
(\sqrt{f_{\mathrm{den}}}-\sqrt{g_{\mathrm{den}}})^2
\le
\frac{(f_{\mathrm{den}}-g_{\mathrm{den}})^2}{4t_{\mathrm{floor}}}.
\]
Apply this with the preceding component floor and density difference, and integrate against the probability reference. We obtain
\[
\begin{aligned}
H_{\mathcal V}^2(\mathbf x)
&\le
\frac{
\bigl(C_\Delta(3/2)^{p_{\mathrm{cmp}}}
p_{\mathrm{cmp}}^2ab\bigr)^2
}{
4(1/2)^{p_{\mathrm{cmp}}}
}\\
&=
\frac{C_\Delta^2}{4}(9/2)^{p_{\mathrm{cmp}}}
p_{\mathrm{cmp}}^4a^2b^2\\
&\le
K_{\mathrm{info}}(9/2)^{p_{\mathrm{cmp}}}
p_{\mathrm{cmp}}^4a^2b^2.
\end{aligned}
\]
% lean: CausalSmith.Stat.TwosamplePointcateAnnotationFrontier.marked_component_hellinger_bound

\item \textit{Conditional factorization and the count bound.}
For fixed covariates, define the augmented graph
\[
\mathscr G_{\mathbf x}
:=
\left(
\{1,\ldots,N\},
\left\{\{i,j\}:
i\ne j,\ x_i,x_j\in C_h,\
\|x_i-x_j\|_\infty\le2\delta
\right\}
\right),
\]
and let
\[
\mathscr K_{\mathbf x}
:=
\{\mathcal V:\mathcal V
\text{ is the vertex set of a connected component of }
\mathscr G_{\mathbf x}\}.
\]
The cube \(C_h\) is that of \cref{obj:synth_3}. Records outside \(C_h\) are isolated in this augmented graph and have likelihood one.

If \(\phi_{\mathbf z}(x_i)\ne0\), then
\[
x_i\in C_h,
\qquad
\left|\frac{x_{i,j_{\mathrm{coord}}}-x_{0,j_{\mathrm{coord}}}}{\delta}-z_{j_{\mathrm{coord}}}\right|<1
\quad(j_{\mathrm{coord}}\in\{1,\ldots,d\}).
\]
Two records sharing a nonzero frame entry therefore satisfy
\[
\|x_i-x_j\|_\infty\le2\delta,
\]
and belong to the same component. Assign each frame index active at any record to this unique component; inactive indices contribute a factor one. The product prior is independent across indices, and each component likelihood depends solely on its assigned indices. Regrouping the finite sign sums by these disjoint groups consequently proves
\[
g_\theta(\mathbf x,\mathbf s)
:=
\mathbb E_\theta
\prod_{i=1}^{N}R_{\theta,\lambda,\eta,i}(\mathbf x,\mathbf s)
=
\prod_{\mathcal V\in\mathscr K_{\mathbf x}}
g_{\theta,\mathcal V}(\mathbf x,\mathbf s).
\]
Each component density depends solely on its component spins. Integrating the fixed-sign record product over those fair spins gives one, because every record likelihood is normalized. Averaging signs therefore gives
\[
g_{\theta,\mathcal V}\ge0,
\qquad
\int g_{\theta,\mathcal V}\,d\mu_{\mathrm{spin}}=1.
\]

Define the component affinities by
\[
\mathcal A_{\mathrm{aff},\mathcal V}(\mathbf x)
:=
\int\sqrt{g_{+1,\mathcal V}g_{-1,\mathcal V}}\,
d\mu_{\mathrm{spin}}.
\]
Normalization and \(2\sqrt{f_{\mathrm{den}}g_{\mathrm{den}}}
\le f_{\mathrm{den}}+g_{\mathrm{den}}\) imply
\[
0\le\mathcal A_{\mathrm{aff},\mathcal V}\le1,
\qquad
H_{\mathcal V}^2
=2(1-\mathcal A_{\mathrm{aff},\mathcal V}).
\]
Integration over disjoint spin blocks multiplies the affinities. The elementary finite-product inequality
\[
1-\prod_{\mathcal V\in\mathscr K_{\mathbf x}}
\mathcal A_{\mathrm{aff},\mathcal V}
\le
\sum_{\mathcal V\in\mathscr K_{\mathbf x}}
(1-\mathcal A_{\mathrm{aff},\mathcal V})
\]
follows by induction from the fact that all the affinities lie in \([0,1]\). Hence
\[
\int(\sqrt{g_{+1}}-\sqrt{g_{-1}})^2\,d\mu_{\mathrm{spin}}
\le
\sum_{\mathcal V\in\mathscr K_{\mathbf x}}
H_{\mathcal V}^2(\mathbf x).
\]

Set
\[
c_I:=\min\{c_H,c_U\}.
\]
Its amplitude conditions permit both the component bound and the uninformative-component equalities. A component containing a labeled record and having size less than two has size one, so its distance is zero by singleton cancellation. A component containing no labeled record is auxiliary-only, so its distance is also zero. Define the informative component counts, for every integer \(p_{\mathrm{cmp}}\ge0\), by
\[
C_{p_{\mathrm{cmp}}}^{\mathrm L}(\mathbf x)
:=
\#\left\{
\mathcal V\in\mathscr K_{\mathbf x}:
|\mathcal V|=p_{\mathrm{cmp}},\
\mathcal V\cap\{1,\ldots,n\}\ne\varnothing
\right\}.
\]
Regrouping the remaining component sum by size gives
\[
\int(\sqrt{g_{+1}}-\sqrt{g_{-1}})^2\,d\mu_{\mathrm{spin}}
\le
K_{\mathrm{info}}a^2b^2
\sum_{p_{\mathrm{cmp}}=2}^{N}
(9/2)^{p_{\mathrm{cmp}}}
p_{\mathrm{cmp}}^4C_{p_{\mathrm{cmp}}}^{\mathrm L}(\mathbf x).
\]
The graph \(\mathscr G_{\mathbf x}\) is exactly the graph in the statement: it retains all \(N\) vertices, including the exterior singleton components, whose density factors equal one.
% lean: CausalSmith.Stat.TwosamplePointcateAnnotationFrontier.marked_conditional_hellinger_count_bound

\item \textit{The labeled-root tree series.}
Define the common covariate probability measure by
\[
\nu_{\mathrm{cov}}
:=
\operatorname{Unif}([0,1]^d)^{\otimes N}.
\]
For \(p_{\mathrm{cmp}}\ge2\), a component counted by
\(C_{p_{\mathrm{cmp}}}^{\mathrm L}\) has a labeled root and a spanning tree on its vertex set. For a fixed root, there are
\(\binom{N-1}{p_{\mathrm{cmp}}-1}\) choices of the other vertices. Orient a spanning tree toward the root. Each nonroot vertex chooses a parent among the \(p_{\mathrm{cmp}}\) selected vertices, so the number of valid rooted trees is at most
\(p_{\mathrm{cmp}}^{p_{\mathrm{cmp}}-1}\).

For a fixed valid tree, integrate its leaf covariates successively. Each child is constrained to a maximum-norm box of side \(4\delta\) around its parent; intersection with the unit cube has uniform mass at most \((4\delta)^d\). The root lies in \(C_h\), whose uniform mass is \(h^d\). Thus the probability of this tree event is at most
\[
h^d\bigl((4\delta)^d\bigr)^{p_{\mathrm{cmp}}-1}.
\]
Dropping the remaining component conditions enlarges the event. Summing these tree witnesses over the \(n\) possible labeled roots gives
\[
\int C_{p_{\mathrm{cmp}}}^{\mathrm L}(\mathbf x)\,
d\nu_{\mathrm{cov}}(\mathbf x)
\le
nh^d
\binom{N-1}{p_{\mathrm{cmp}}-1}
p_{\mathrm{cmp}}^{p_{\mathrm{cmp}}-1}
\bigl((4\delta)^d\bigr)^{p_{\mathrm{cmp}}-1}.
\]
For \(p_{\mathrm{cmp}}>N\), the count and the binomial coefficient are both zero.

The binomial bound and the exponential series imply
\[
\binom{N-1}{p_{\mathrm{cmp}}-1}
\le\frac{N^{p_{\mathrm{cmp}}-1}}{(p_{\mathrm{cmp}}-1)!},
\qquad
\frac{p_{\mathrm{cmp}}^{p_{\mathrm{cmp}}-1}}
{(p_{\mathrm{cmp}}-1)!}
\le \exp(p_{\mathrm{cmp}}).
\]
For the second inequality, its left side is one nonnegative summand of the exponential series at \(p_{\mathrm{cmp}}\). Consequently
\[
\binom{N-1}{p_{\mathrm{cmp}}-1}
p_{\mathrm{cmp}}^{p_{\mathrm{cmp}}-1}
\le N^{p_{\mathrm{cmp}}-1}\exp(p_{\mathrm{cmp}}).
\]

Define
\[
\kappa_{\mathrm{tree}}:=\frac92,
\qquad
S_{\mathrm{tree}}
:=
\sum_{j=0}^{\infty}(j+2)^4\,2^{-j},
\qquad
E_{\mathrm{occ}}
:=
\kappa_{\mathrm{tree}}\exp(1)4^d,
\]
\[
c_0:=\frac1{2E_{\mathrm{occ}}},
\qquad
C_{\mathrm{tree}}
:=
(S_{\mathrm{tree}}+1)
\bigl(\kappa_{\mathrm{tree}}\exp(1)\bigr)^2\,4^d.
\]
The series defining \(S_{\mathrm{tree}}\) converges because it is a shifted polynomial-geometric series. These definitions give \(c_0>0\) and \(C_{\mathrm{tree}}>0\).

Under \(N\delta^d\le c_0\), define
\[
z_{\mathrm{occ}}:=E_{\mathrm{occ}}N\delta^d,
\qquad
0\le z_{\mathrm{occ}}\le\frac12.
\]
Applying the binomial-tree bound with \(p_{\mathrm{cmp}}=j+2\) gives
\[
\begin{aligned}
&\kappa_{\mathrm{tree}}^{j+2}(j+2)^4
\binom{N-1}{j+1}(j+2)^{j+1}
\bigl((4\delta)^d\bigr)^{j+1}\\
&\qquad\le
\kappa_{\mathrm{tree}}\exp(1)(j+2)^4
z_{\mathrm{occ}}^{j+1}.
\end{aligned}
\]
For every finite truncation,
\[
\sum_{j=0}^{N-1}(j+2)^4z_{\mathrm{occ}}^{j+1}
\le
z_{\mathrm{occ}}\sum_{j=0}^{N-1}(j+2)^4\,2^{-j}
\le z_{\mathrm{occ}}S_{\mathrm{tree}}.
\]
Multiplying the tree-count inequality by its size weight and summing therefore yields
\[
\begin{aligned}
&\sum_{j=0}^{N-1}
\kappa_{\mathrm{tree}}^{j+2}(j+2)^4
\int C_{j+2}^{\mathrm L}\,d\nu_{\mathrm{cov}}\\
&\qquad\le
nh^d\kappa_{\mathrm{tree}}\exp(1)
z_{\mathrm{occ}}S_{\mathrm{tree}}\\
&\qquad\le C_{\mathrm{tree}}nNh^d\delta^d.
\end{aligned}
\]
The final inequality replaces \(S_{\mathrm{tree}}\) by
\(S_{\mathrm{tree}}+1\) and substitutes the definition of
\(z_{\mathrm{occ}}\).
% lean: CausalSmith.Stat.TwosamplePointcateAnnotationFrontier.labeled_component_series_bound

\item \textit{Integration for the original experiment.}
The covariate distribution is \(\nu_{\mathrm{cov}}\) for every sign realization and hypothesis. Conditional on these covariates, the original records have density \(g_\theta\) against \(\mu_{\mathrm{spin}}\). Thus, after grouping covariates and spins, the mixture laws have densities \(g_{+1}\) and \(g_{-1}\) against
\(\nu_{\mathrm{cov}}\otimes\mu_{\mathrm{spin}}\).

For two densities against a common dominating measure, substitution into \cref{obj:def:hellinger} gives the integral of their squared square-root difference. Applying this identity and then integrating over covariates yields
\[
H^2(M_{+1},M_{-1})
=
\int
\left[
\int(\sqrt{g_{+1}(\mathbf x,\mathbf s)}
-\sqrt{g_{-1}(\mathbf x,\mathbf s)})^2
\,d\mu_{\mathrm{spin}}(\mathbf s)
\right]d\nu_{\mathrm{cov}}(\mathbf x).
\]
The finite graph counts are measurable, and all functions in the finite weighted count sum are integrable. Integrating the conditional bound therefore gives
\[
H^2(M_{+1},M_{-1})
\le
K_{\mathrm{info}}a^2b^2
\sum_{p_{\mathrm{cmp}}=2}^{N}
(9/2)^{p_{\mathrm{cmp}}}p_{\mathrm{cmp}}^4
\int C_{p_{\mathrm{cmp}}}^{\mathrm L}\,d\nu_{\mathrm{cov}}.
\]
Every term is nonnegative, so extending the size sum to the shifted range used above gives
\[
\begin{aligned}
H^2(M_{+1},M_{-1})
&\le
K_{\mathrm{info}}a^2b^2
\sum_{j=0}^{N-1}
(9/2)^{j+2}(j+2)^4
\int C_{j+2}^{\mathrm L}\,d\nu_{\mathrm{cov}}\\
&\le
K_{\mathrm{info}}C_{\mathrm{tree}}
nNh^d\delta^da^2b^2.
\end{aligned}
\]
Define
\[
C:=K_{\mathrm{info}}C_{\mathrm{tree}}>0.
\]
This proves the required information estimate under the amplitude conditions with \(c_I\) and the occupancy condition with \(c_0\).
% lean: CausalSmith.Stat.TwosamplePointcateAnnotationFrontier.marked_information_bound

\item \textit{Common calibration and the remaining equalities.}
Choose
\[
c:=\min\{c_M,\min\{c_S,\min\{c_U,c_I\}\}\}.
\]
All four constants are positive and at most one, so
\[
0<c\le1,
\qquad
c\le c_M,\quad c\le c_S,\quad c\le c_U,\quad c\le c_I.
\]
Fix a tuple satisfying the theorem's hypotheses with this \(c\). Since \(\delta>0\) and \(h>0\), multiplication by their real powers preserves these comparisons. For each
\(c'\in\{c_M,c_S,c_U,c_I\}\),
\[
a\le c\delta^\alpha\le c'\delta^\alpha,\qquad
b\le c\delta^\beta\le c'\delta^\beta,\qquad
ab\le ch^\gamma\le c'h^\gamma.
\]
The membership, singleton, uninformative-component, and information estimates established above therefore apply simultaneously. Weight normalization makes the finite pushforward priors probability measures, and the pointwise membership conclusion supports them on \(\mathcal M\).

For the propensity distribution, the sign flip preserves the covariate-treatment marginal:
\[
(P_{+1,\lambda,\eta})_{XA}
=
(P_{-1,\lambda,-\eta})_{XA}.
\]
The designated propensity witnesses of these admitted laws are continuous and are versions of the same conditional treatment probability. Uniqueness of conditional margins followed by continuity under the full-support uniform design gives
\[
e_{P_{+1,\lambda,\eta}}(x)
=
e_{P_{-1,\lambda,-\eta}}(x)
\qquad(x\in[0,1]^d).
\]
Reindexing the finite pushforwards using
\(w_{+1}(\lambda,\eta)=w_{-1}(\lambda,-\eta)\) proves equality of the distributions of the entire propensity functions restricted to the cube.

At the target, \(B_h(x_0)=1\), and hence
\[
\tau_\theta(x_0)=-2\theta ab,
\qquad
\tau_{-1}(x_0)-\tau_{+1}(x_0)=4ab.
\]
Because \(a,b>0\) and \(c\le1\le4\),
\[
cab=c(ab)\le4(ab)=4ab.
\]

Finally, for every integer \(k_{\mathrm{aux}}\ge0\), reindexing by the same sign flip gives
\[
\begin{aligned}
\int P_{XA}^{\otimes k_{\mathrm{aux}}}\,d\varpi_{+1}(P)
&=
\sum_{\lambda,\eta}
w_{+1}(\lambda,\eta)
(P_{+1,\lambda,\eta})_{XA}^{\otimes k_{\mathrm{aux}}}\\
&=
\sum_{\lambda,\eta}
w_{-1}(\lambda,\eta)
(P_{-1,\lambda,\eta})_{XA}^{\otimes k_{\mathrm{aux}}}\\
&=
\int P_{XA}^{\otimes k_{\mathrm{aux}}}\,d\varpi_{-1}(P).
\end{aligned}
\]
For any fixed auxiliary covariates, their conditional mixture kernels are
\[
\sum_{\lambda,\eta}
w_\theta(\lambda,\eta)
\bigotimes_{i=1}^{k_{\mathrm{aux}}}
\operatorname{Bernoulli}(e_\lambda(x_i)).
\]
The factors are unchanged by the sign flip, so these kernels also agree for every covariate collection. When \(k_{\mathrm{aux}}=0\), the products are the unit probability measure on the empty tuple, and normalization gives the same equality.

The singleton calculation supplies the stated uniform labeled conditional laws. The conditional factorization and component equalities apply to every cube covariate collection and every spin assignment. Together with the calibrated information bound, these establish all assertions.
% lean: CausalSmith.Stat.TwosamplePointcateAnnotationFrontier.marked_component_certificate
\end{enumerate}
\end{proof}

\begin{proof}[Proof of \cref{obj:prop:annotation-threshold}]
Fix the public parameters in the stated domain. The frontier, oracle, and acquisition constants below are finite real numbers depending only on these parameters. In particular,
\[
\gamma>0,\qquad S>0,\qquad 2\gamma+d>0,\qquad \Delta>0.
\]
For an arbitrary auxiliary-size sequence, write
\[
N_n:=n+m_n\qquad(n\in\mathbb N).
\]
All asymptotic arguments may be restricted to the tail \(n\ge2\).

\begin{enumerate}
\item\label{proof:annotation-threshold-risk-rate}
By \cref{obj:thm:sharp-annotation-frontier,obj:thm:supplied-propensity}, choose constants satisfying
\[
0<c_F\le C_F<\infty,\qquad 0<c_O\le C_O<\infty
\]
such that, uniformly for \(n\ge2\) and \(m\ge0\),
\[
c_F r_*(n,m)\le R(n,m)\le C_F r_*(n,m),
\qquad
c_O r_o(n)\le R^e(n,m)\le C_O r_o(n).
\]
Since
\[
r_*(n,m)\ge r_o(n)=n^{-\gamma/(2\gamma+d)}>0,
\]
these bounds make both risks positive and finite. Dividing the upper bound for one risk by the lower bound for the other gives
\[
\begin{aligned}
0\le \frac{R(n,m)}{R^e(n,m)}
&\le \frac{C_F}{c_O}\frac{r_*(n,m)}{r_o(n)},\\
0\le \frac{r_*(n,m)}{r_o(n)}
&\le \frac{C_O}{c_F}\frac{R(n,m)}{R^e(n,m)}.
\end{aligned}
\]
The same calculation, using the frontier bounds at \((n,m)\) and \((n,0)\), gives
\[
\begin{aligned}
0\le \frac{R(n,m)}{R(n,0)}
&\le \frac{C_F}{c_F}\frac{r_*(n,m)}{r_*(n,0)},\\
0\le \frac{r_*(n,m)}{r_*(n,0)}
&\le \frac{C_F}{c_F}\frac{R(n,m)}{R(n,0)}.
\end{aligned}
\]
The first pair transfers eventual boundedness in both directions. The second pair transfers convergence to zero by squeezing between zero and a fixed multiple of the other ratio. Moreover, positivity of \(r_o(n)\) gives
\[
\frac{r_*(n,m_n)}{r_o(n)}=O(1)
\quad\Longleftrightarrow\quad
r_*(n,m_n)=O(r_o(n)).
\]
Consequently,
\[
\begin{aligned}
\frac{R(n,m_n)}{R^e(n,m_n)}=O(1)
&\quad\Longleftrightarrow\quad r_*(n,m_n)=O(r_o(n)),\\
\frac{R(n,m_n)}{R(n,0)}\longrightarrow0
&\quad\Longleftrightarrow\quad
\frac{r_*(n,m_n)}{r_*(n,0)}\longrightarrow0.
\end{aligned}
\]

If the rate comparison holds, there is a real constant \(K_{\mathrm{rate}}\) such that
\[
r_*(n,m_n)\le K_{\mathrm{rate}}r_o(n)
\qquad\text{eventually}.
\]
Set
\[
\zeta:=\max\{K_{\mathrm{rate}},1\}.
\]
Then \(1\le\zeta<\infty\), and the defining inequality in
\cref{obj:def:annotation-region} holds eventually. Conversely, eventual membership in that region supplies the rate comparison with constant \(\zeta\). Thus
\[
r_*(n,m_n)=O(r_o(n))
\quad\Longleftrightarrow\quad
\exists\,\zeta\in[1,\infty)\ 
\text{such that }(n,m_n)\in\mathcal B_*(\zeta)
\text{ eventually}.
\]
% lean: CausalSmith.Stat.TwosamplePointcateAnnotationFrontier.annotation_risk_rate_equivalences

\item
For the extended nonnegative lower limit, the tail characterization is
\[
0<\liminf_{n\to\infty}\frac{N_n}{n^{q_*}}
\quad\Longleftrightarrow\quad
\exists\,c_{\mathrm{count}}>0\
\text{such that }
c_{\mathrm{count}}\le\frac{N_n}{n^{q_*}}
\text{ eventually}.
\]
Indeed, a positive lower limit admits a smaller positive finite real number, which bounds every sufficiently late term from below; this also applies when the lower limit is \(+\infty\). Conversely, an eventual positive lower bound bounds the lower limit from below.

We now translate annotation-region membership into precisely such a bound. Directly from the definitions,
\[
q_*=\frac{S_{\mathrm{crit}}}{S},
\qquad
\Delta=(2\gamma+d)(1+q_*).
\]
For \(n\ge2\), this yields
\[
r_o(n)^{-\Delta/\gamma}
=n^{\Delta/(2\gamma+d)}
=n^{1+q_*}
=n\,n^{q_*}.
\]
Let \(\zeta\ge1\). Since \(r_o(n)\le\zeta r_o(n)\), the maximum defining \(r_*\) and
\cref{obj:def:annotation-region} give
\[
\begin{aligned}
(n,m)\in\mathcal B_*(\zeta)
&\quad\Longleftrightarrow\quad
(nN)^{-\gamma/\Delta}\le\zeta r_o(n)\\
&\quad\Longleftrightarrow\quad
(\zeta r_o(n))^{-\Delta/\gamma}\le nN\\
&\quad\Longleftrightarrow\quad
\zeta^{-\Delta/\gamma}\le\frac{N}{n^{q_*}}.
\end{aligned}
\]
Here the second equivalence uses the strictly decreasing positive-power map with exponent \(-\gamma/\Delta\), whose inverse has exponent \(-\Delta/\gamma\); all bases are positive.

Eventual membership therefore supplies the positive lower bound
\[
c_{\mathrm{count}}:=\zeta^{-\Delta/\gamma}>0.
\]
In the reverse direction, suppose an eventual lower bound \(c_{\mathrm{count}}>0\) is given, and choose
\[
\zeta:=\max\left\{1,c_{\mathrm{count}}^{-\gamma/\Delta}\right\}.
\]
Then
\[
1\le\zeta<\infty,
\qquad
\zeta^{-\Delta/\gamma}
\le
\left(c_{\mathrm{count}}^{-\gamma/\Delta}\right)^{-\Delta/\gamma}
=c_{\mathrm{count}}.
\]
The displayed count characterization gives eventual membership in
\(\mathcal B_*(\zeta)\). Combining this with \cref{proof:annotation-threshold-risk-rate} proves all four oracle-order equivalences.
% lean: CausalSmith.Stat.TwosamplePointcateAnnotationFrontier.annotationRegion_liminf_iff

\item
To characterize strict improvement, first suppose \(S<S_{\mathrm{crit}}\). The identities above imply
\[
q_*=\frac{S_{\mathrm{crit}}}{S}>1,
\qquad
\Delta=(2\gamma+d)(1+q_*)>2(2\gamma+d),
\]
and hence
\[
\frac{2\gamma}{\Delta}<\frac{\gamma}{2\gamma+d}.
\]
Define the positive decay exponent
\[
\eta_{\mathrm{dec}}
:=\frac{\gamma}{2\gamma+d}-\frac{2\gamma}{\Delta}>0.
\]
For \(n\ge2\), the inequality \(q_*>1\) gives
\[
n=n^1\le n^{q_*}.
\]
The lower-count branch in
\cref{obj:thm:sharp-annotation-frontier}, applied with \(m=0\), therefore gives
\[
r_*(n,0)=(n^2)^{-\gamma/\Delta}
=n^{-2\gamma/\Delta}.
\]
Dividing the two terms in the maximum defining \(r_*(n,m_n)\) by this positive rate yields
\[
\frac{r_o(n)}{n^{-2\gamma/\Delta}}
=n^{-\eta_{\mathrm{dec}}},
\qquad
\frac{(nN_n)^{-\gamma/\Delta}}{n^{-2\gamma/\Delta}}
=\left(\frac{N_n}{n}\right)^{-\gamma/\Delta}.
\]
Thus
\[
\frac{r_*(n,m_n)}{r_*(n,0)}
=
\max\left\{
n^{-\eta_{\mathrm{dec}}},
\left(\frac{N_n}{n}\right)^{-\gamma/\Delta}
\right\}.
\]

The second term tends to zero exactly when \(N_n/n\to+\infty\). To see the reverse implication explicitly, let
\[
b_{\mathrm{thr}}\in\mathbb R,
\qquad
x_{\mathrm{thr}}:=\max\{1,b_{\mathrm{thr}}\}>0.
\]
If the second term tends to zero, then eventually
\[
\left(\frac{N_n}{n}\right)^{-\gamma/\Delta}
\le x_{\mathrm{thr}}^{-\gamma/\Delta}.
\]
Inverting this decreasing power gives
\[
\frac{N_n}{n}
\ge
\left(x_{\mathrm{thr}}^{-\gamma/\Delta}\right)^{-\Delta/\gamma}
=x_{\mathrm{thr}}
\ge b_{\mathrm{thr}}.
\]
This proves divergence to \(+\infty\). The forward implication follows from the convergence of a strictly negative power to zero at \(+\infty\). Since
\[
\frac{N_n}{n}=1+\frac{m_n}{n}\qquad(n\ge2),
\]
divergence of the total-to-labeled ratio is equivalent to divergence of the auxiliary-to-labeled ratio: an eventual lower bound at any real threshold for either ratio gives the corresponding bound for the other after shifting the threshold by one.

The first term \(n^{-\eta_{\mathrm{dec}}}\) tends to zero. Both terms are nonnegative, so convergence of their maximum to zero implies convergence of the second term by squeezing; conversely, convergence of both terms implies convergence of their maximum. We have proved, in this regime,
\[
\frac{r_*(n,m_n)}{r_*(n,0)}\longrightarrow0
\quad\Longleftrightarrow\quad
\frac{m_n}{n}\longrightarrow+\infty.
\]

For the complementary case \(S\ge S_{\mathrm{crit}}\),
\cref{obj:thm:sharp-annotation-frontier} gives, for every \(n\ge2\),
\[
r_*(n,m_n)=r_*(n,0)=r_o(n)>0,
\qquad
\frac{r_*(n,m_n)}{r_*(n,0)}=1.
\]
The ratio therefore has limit one, which excludes convergence to zero. Combining the two cases and the risk-to-rate equivalence proves
\[
\frac{R(n,m_n)}{R(n,0)}\longrightarrow0
\quad\Longleftrightarrow\quad
\frac{r_*(n,m_n)}{r_*(n,0)}\longrightarrow0
\quad\Longleftrightarrow\quad
S<S_{\mathrm{crit}}
\ \text{and}\ 
\frac{m_n}{n}\longrightarrow+\infty.
\]
% lean: CausalSmith.Stat.TwosamplePointcateAnnotationFrontier.annotation_rate_improvement_iff

\item
The claimed exponent comparison follows from the same positive-denominator identity:
\[
S<S_{\mathrm{crit}}
\quad\Longrightarrow\quad
q_*=\frac{S_{\mathrm{crit}}}{S}>1,
\qquad S>0.
\]
% lean: CausalSmith.Stat.TwosamplePointcateAnnotationFrontier.qStar_gt_one

\item
Suppose \(m_n/n=O(1)\). Choose a real constant \(K_{\mathrm{aux}}\) satisfying
\[
\frac{m_n}{n}\le K_{\mathrm{aux}}
\qquad\text{eventually},
\]
and define
\[
K_{\mathrm{aux}}^+:=\max\{K_{\mathrm{aux}},0\},
\qquad
M_{\mathrm{aux}}
:=(1+K_{\mathrm{aux}}^+)^{\gamma/\Delta}\ge1.
\]
Eventually \(m_n\le K_{\mathrm{aux}}^+n\), and consequently
\[
n^2\le nN_n\le(1+K_{\mathrm{aux}}^+)n^2.
\]
Taking the negative power \(-\gamma/\Delta\) gives
\[
(nN_n)^{-\gamma/\Delta}\le(n^2)^{-\gamma/\Delta},
\]
and
\[
M_{\mathrm{aux}}^{-1}(n^2)^{-\gamma/\Delta}
\le(nN_n)^{-\gamma/\Delta},
\qquad
(n^2)^{-\gamma/\Delta}
\le M_{\mathrm{aux}}(nN_n)^{-\gamma/\Delta}.
\]
The oracle term is common to both rates and satisfies
\[
r_o(n)\le M_{\mathrm{aux}}r_o(n).
\]
Taking maxima therefore gives
\[
r_*(n,m_n)\le r_*(n,0)
\le M_{\mathrm{aux}}r_*(n,m_n)
\qquad\text{eventually}.
\]
Using the risk sandwiches,
\[
\begin{aligned}
R(n,m_n)
&\le C_F r_*(n,m_n)
\le C_F r_*(n,0)
\le \frac{C_F}{c_F}R(n,0),\\
R(n,0)
&\le C_F r_*(n,0)
\le C_F M_{\mathrm{aux}}r_*(n,m_n)
\le \frac{C_F M_{\mathrm{aux}}}{c_F}R(n,m_n).
\end{aligned}
\]
These eventual inequalities prove both stated big-\(O\) comparisons.
% lean: CausalSmith.Stat.TwosamplePointcateAnnotationFrontier.bounded_auxiliary_risk_order

\item
If \(S\ge S_{\mathrm{crit}}\), the frontier rate equals \(r_o(n)\). For every \(n\ge2\) and \(m\ge0\), the frontier upper bound and the oracle lower bound consequently give
\[
R(n,m)
\le C_F r_o(n)
=\frac{C_F}{c_O}\bigl(c_O r_o(n)\bigr)
\le\frac{C_F}{c_O}R^e(n,m).
\]
Dividing by the positive oracle risk yields
\[
0\le\frac{R(n,m)}{R^e(n,m)}
\le\frac{C_F}{c_O}.
\]
This uniform bound proves the oracle-order comparison for every auxiliary-size sequence.
% lean: CausalSmith.Stat.TwosamplePointcateAnnotationFrontier.smooth_oracle_risk_ratio

\item
For the finite-sample conclusions,
\cref{obj:thm:sharp-annotation-frontier} supplies, simultaneously for all \(n\ge2\) and \(m\ge0\),
\[
c_F r_*(n,m)\le R(n,m)
\le\sup_{P\in\mathcal M}\mathcal R_{n,m}(T_*,P)
\le C_F r_*(n,m).
\]
In particular, its upper bound follows by comparing the minimax value with the worst-case risk of \(T_*\). The same cited statement supplies exactly
\[
r_*(n,m)=
\begin{cases}
r_o(n),&S\ge S_{\mathrm{crit}},\\
(nN)^{-\gamma/\Delta},
&S<S_{\mathrm{crit}},\quad N\le n^{q_*},\\
r_o(n),
&S<S_{\mathrm{crit}},\quad N\ge n^{q_*},
\end{cases}
\]
together with
\[
N=n^{q_*}
\quad\Longrightarrow\quad
(nN)^{-\gamma/\Delta}=r_o(n).
\]
Thus \(c_F,C_F\) furnish the asserted finite-sample constants, including the equality of the two count branches at their boundary.
% lean: CausalSmith.Stat.TwosamplePointcateAnnotationFrontier.sharp_annotation_frontier

\item
For \(n\ge2\), specialize to \(m=0\). If \(S<S_{\mathrm{crit}}\), then \(q_*>1\), so
\[
N=n\le n^{q_*},
\qquad
r_*(n,0)=(n^2)^{-\gamma/\Delta}
=n^{-2\gamma/\Delta}.
\]
If \(S\ge S_{\mathrm{crit}}\), the first branch of
\cref{obj:thm:sharp-annotation-frontier} gives \(r_*(n,0)=r_o(n)\). Hence
\[
r_*(n,0)=
\begin{cases}
n^{-2\gamma/\Delta},&S<S_{\mathrm{crit}},\\
r_o(n),&S\ge S_{\mathrm{crit}}.
\end{cases}
\]
% lean: CausalSmith.Stat.TwosamplePointcateAnnotationFrontier.supervised_rate_branches

\item
Define the average nuisance smoothness by
\[
s_{\mathrm{pub}}:=\frac S2>0.
\]
By \cref{obj:lem:published-cate-benchmark}, algebra with positive denominators gives
\[
\frac{d/4}{1+d/(2\gamma)}
=\frac{S_{\mathrm{crit}}}{2},
\qquad
\frac{1}{1+d/(2\gamma)+d/(4s_{\mathrm{pub}})}
=\frac{2\gamma}{\Delta},
\qquad
\frac{1}{2+d/\gamma}
=\frac{\gamma}{2\gamma+d}.
\]
In particular,
\[
s_{\mathrm{pub}}<\frac{d/4}{1+d/(2\gamma)}
\quad\Longleftrightarrow\quad
S<S_{\mathrm{crit}}.
\]
For \(n\ge2\), these identities identify the numerical benchmark's strict split and its two exponents with the supervised branches just established.

The remaining sample sizes are \(n=0\) and \(n=1\). At \(n=1\), all power bases equal one, so both the benchmark and \(r_*(1,0)\) equal one. At zero, the numerical formulas use the convention
\[
0^{t_{\mathrm{pow}}}=0
\qquad
\bigl(t_{\mathrm{pow}}\in\mathbb R\setminus\{0\}\bigr).
\]
All exponents occurring here are strictly negative:
\[
-\frac{\gamma}{2\gamma+d}<0,
\qquad
-\frac{\gamma}{\Delta}<0,
\qquad
-\frac{2\gamma}{\Delta}<0.
\]
Thus both branches of the benchmark equal zero at \(n=0\), and the two terms in the maximum defining \(r_*(0,0)\) also equal zero. This proves the benchmark equality for every \(n\in\mathbb N\), with the stated formulas for \(n\ge1\).
% lean: CausalSmith.Stat.TwosamplePointcateAnnotationFrontier.published_benchmark_eq_supervised_rate

\item
For acquisition, define the positive exponents
\[
p_0:=\frac{2\gamma+d}{\gamma}
=2+\frac d\gamma,
\qquad
p_1:=\frac{\Delta}{\gamma}
=2+\frac d\gamma+\frac dS.
\]
For positive \(\eta\), the power-scaling identity is
\[
\left(\frac{\eta}{c_{\mathrm{scale}}}\right)^{-p_{\mathrm{scale}}}
=
c_{\mathrm{scale}}^{p_{\mathrm{scale}}}
\eta^{-p_{\mathrm{scale}}}
\qquad
\left(
c_{\mathrm{scale}}\in\{c_F,C_F\},\
p_{\mathrm{scale}}\in\{p_0,p_1\}
\right).
\]
Choose
\[
c_{\mathrm{acq}}
:=\min\{c_F^{p_0},c_F^{p_1}\}>0,
\qquad
C_{\mathrm{acq}}
:=\max\{C_F^{p_0},C_F^{p_1}\}>0.
\]

For any positive \(\eta\), \(n\ge2\), and \(m\ge0\), the defining maximum in
\cref{obj:def:sharp-rate} gives
\[
\begin{aligned}
r_*(n,m)\le\eta
&\quad\Longleftrightarrow\quad
n^{-\gamma/(2\gamma+d)}\le\eta
\ \text{and}\
(nN)^{-\gamma/\Delta}\le\eta\\
&\quad\Longleftrightarrow\quad
\eta^{-p_0}\le n
\ \text{and}\
\eta^{-p_1}\le nN.
\end{aligned}
\]
The second equivalence again inverts strictly decreasing powers of positive bases.

Now let \(0<\eta<1\), and suppose a Borel decision \(T\) satisfies
\[
\mathcal R_{n,m}(T,P)\le\eta
\qquad(P\in\mathcal M).
\]
By \cref{obj:def:minimax} and the frontier lower bound,
\[
c_F r_*(n,m)\le R(n,m)
\le\sup_{P\in\mathcal M}\mathcal R_{n,m}(T,P)
\le\eta.
\]
Therefore \(r_*(n,m)\le\eta/c_F\). Apply the exact rate equivalence at the positive scale \(\eta/c_F\), and then the scaling identity, to obtain
\[
c_F^{p_0}\eta^{-p_0}\le n,
\qquad
c_F^{p_1}\eta^{-p_1}\le nN.
\]
The choice of the minimum gives the common necessity constant:
\[
c_{\mathrm{acq}}\eta^{-p_0}
\le c_F^{p_0}\eta^{-p_0}\le n,
\qquad
c_{\mathrm{acq}}\eta^{-p_1}
\le c_F^{p_1}\eta^{-p_1}\le nN.
\]

Conversely, suppose
\[
C_{\mathrm{acq}}\eta^{-p_0}\le n,
\qquad
C_{\mathrm{acq}}\eta^{-p_1}\le nN.
\]
By the choice of the maximum and the scaling identity,
\[
\left(\frac{\eta}{C_F}\right)^{-p_0}
=C_F^{p_0}\eta^{-p_0}
\le C_{\mathrm{acq}}\eta^{-p_0}\le n,
\]
and
\[
\left(\frac{\eta}{C_F}\right)^{-p_1}
=C_F^{p_1}\eta^{-p_1}
\le C_{\mathrm{acq}}\eta^{-p_1}\le nN.
\]
The exact equivalence at scale \(\eta/C_F\) gives
\[
r_*(n,m)\le\frac{\eta}{C_F}.
\]
By \cref{obj:thm:sharp-annotation-frontier}, the decision \(T_*\) from
\cref{obj:def:sharp-decision} is Borel measurable, and, for every \(P\in\mathcal M\),
\[
\mathcal R_{n,m}(T_*,P)
\le\sup_{P'\in\mathcal M}\mathcal R_{n,m}(T_*,P')
\le C_F r_*(n,m)
\le\eta.
\]
Thus \(c_{\mathrm{acq}},C_{\mathrm{acq}}\) supply the asserted acquisition constants. Finally,
\[
n\le n+m=N
\]
because \(m\ge0\).
% lean: CausalSmith.Stat.TwosamplePointcateAnnotationFrontier.annotation_acquisition
\end{enumerate}
\end{proof}
